% ROZDZIAL 4 -- Variadic templates w C++ (C++20)
% Fragment dolaczany do glownego dokumentu po rozdziale 3.
% Wymagane w preambule: \usepackage{listings}, konfiguracja lstlisting
% (taka sama jak w rozdzialach 2 i 3).
% Uzycie: % ROZDZIAL 4 -- Variadic templates w C++ (C++20)
% Fragment dolaczany do glownego dokumentu po rozdziale 3.
% Wymagane w preambule: \usepackage{listings}, konfiguracja lstlisting
% (taka sama jak w rozdzialach 2 i 3).
% Uzycie: % ROZDZIAL 4 -- Variadic templates w C++ (C++20)
% Fragment dolaczany do glownego dokumentu po rozdziale 3.
% Wymagane w preambule: \usepackage{listings}, konfiguracja lstlisting
% (taka sama jak w rozdzialach 2 i 3).
% Uzycie: % ROZDZIAL 4 -- Variadic templates w C++ (C++20)
% Fragment dolaczany do glownego dokumentu po rozdziale 3.
% Wymagane w preambule: \usepackage{listings}, konfiguracja lstlisting
% (taka sama jak w rozdzialach 2 i 3).
% Uzycie: \input{rozdzial_4_variadic_templates_cpp.tex}

\chapter[Variadic templates w C++]{Szablony o zmiennej liczbie argumentów (\textit{variadic templates})}
\label{chap:variadic-templates-cpp}

\section{Idea i zastosowania}

Zwykły szablon funkcji przyjmuje z góry ustaloną liczbę parametrów
szablonu. \emph{Variadic template} (szablon o zmiennej liczbie parametrów)
może natomiast przyjmować \emph{dowolną liczbę} argumentów szablonu,
również zero. Mechanizmem, który to umożliwia, jest
\emph{pakiet parametrów} (ang. \emph{parameter pack}).

\begin{lstlisting}
template<class T>
void single(const T&) {}        // one type

template<class... Ts>
void many(const Ts&...) {}      // zero or more types

void demo() {
    many();                     // zero arguments
    many(10);                   // one argument
    many(10, 3.5, "abc");      // three arguments
}
\end{lstlisting}

Pakiety parametrów są ważną częścią nowoczesnego C++.
Wykorzystują je m.in. klasy \texttt{std::tuple} i \texttt{std::variant}
oraz fabryki \texttt{std::make\_unique} i \texttt{std::make\_shared}.
Nie trzeba dzięki nim pisać przeciążeń dla każdej
liczby argumentów.

\paragraph{Wymagane wersje języka.}
Pakiety parametrów i perfect forwarding pojawiły się w C++11.
\emph{Fold expressions} dodano w C++17, a koncepcje
(\texttt{concepts}) i klauzulę \texttt{requires} w C++20.
Dalej zakładamy C++20.

\section{Trzy kropki: deklaracja pakietu a rozwinięcie}

\begin{center}\small
\begin{tabularx}{\textwidth}{@{}p{0.35\textwidth}X@{}}
\toprule
Zapis & Działanie \\
\midrule
\code{typename... Ts} & Deklaruje pakiet parametrów typowych szablonu. \\
\code{Ts...} & Rozwija pakiet typów w miejscu, które dopuszcza taki wzorzec. \\
\code{args...} & Rozwija pakiet argumentów, tworząc listę wyrażeń. \\
\code{sizeof...(Ts)} & Podaje liczbę elementów pakietu w czasie kompilacji. \\
\code{std::forward<Ts>(args)...} & Przekazuje każdy argument z zachowaniem jego kategorii wartości. \\
\bottomrule
\end{tabularx}
\end{center}
\label{sec:variadic-trzy-kropki}

W różnych miejscach token \texttt{...} oznacza różne operacje.
Rozróżnienie to stanowi podstawę poprawnego używania szablonów
wariadycznych.

\begin{lstlisting}
template<class... Ts>             // declare type pack Ts
void f(Ts... xs) {                // declare function parameter pack xs
    // ...
}
\end{lstlisting}

Deklaracja \texttt{class... Ts} wprowadza pakiet \emph{typów}.
Zapis \texttt{Ts... xs} rozwija pakiet typów i wprowadza odpowiadający
mu pakiet parametrów funkcji. Wywołanie
\texttt{f(1, 2.0, 'x')} tworzy w uproszczeniu funkcję o parametrach
\texttt{(int, double, char)}.

\subsection{Rozwinięcie pakietu (\textit{pack expansion})}

Rozwinięcie składa się ze \emph{wzorca} (\emph{pattern}) oraz następującego
po nim \texttt{...}. Dla każdego elementu pakietu kompilator podstawia
kolejny element do całego wzorca:

\begin{lstlisting}
template<class... Ts>
void demo(Ts... xs) {
    g(xs...);                // g(x1, x2, ..., xn)
    h((xs * 2)...);          // h(x1*2, x2*2, ..., xn*2)
}
\end{lstlisting}

W drugim wywołaniu wzorcem jest \texttt{(xs * 2)}, a \emph{nie}
pojedyncza zmienna \texttt{xs}. Ważne: \texttt{xs...} nie jest
kontenerem ani tablicą; nie można po nim bezpośrednio iterować zwykłą
pętlą \texttt{for}. To składnia kompilatora służąca do generowania
listy wyrażeń.

\paragraph{Pakiet typów a pakiet wartości.}
W przykładzie \texttt{Ts} reprezentuje typy, natomiast
\texttt{xs} to konkretne parametry funkcji. Jeden typ może wystąpić
w pakiecie wielokrotnie, a każdy element pakietu wartości zachowuje
swój właściwy typ.

\subsection{Liczba elementów: \texttt{sizeof...}}

Operator \texttt{sizeof...} daje liczbę elementów pakietu jako
stałą czasu kompilacji. Nie należy go mylić z operatorem
\texttt{sizeof} mierzącym rozmiar typu lub obiektu.

\begin{lstlisting}
#include <cstddef>

template<class... Ts>
constexpr std::size_t count(Ts&&... xs) {
    static_assert(sizeof...(Ts) == sizeof...(xs));
    return sizeof...(xs);
}

static_assert(count() == 0);
static_assert(count(2, 3.5, 'a') == 3);
\end{lstlisting}

Nazwy \texttt{xs} nie używamy w ciele funkcji do odczytywania
wartości; operacja liczenia jest czysto składniowa.

\section{Pierwsze praktyczne funkcje}

\subsection{Wypisywanie dowolnej liczby argumentów}

Najprostszy sposób wykorzystania pakietu to jego rozwinięcie
w wywołaniu innej funkcji:

\begin{lstlisting}
#include <iostream>
#include <string>

void output(int x, double y, const std::string& z) {
    std::cout << x << ' ' << y << ' ' << z << '\n';
}

template<class... Args>
void forward_to_output(Args... args) {
    output(args...);
}

int main() {
    forward_to_output(3, 2.5, std::string{"hello"});
}
\end{lstlisting}

Funkcja jest poprawna tylko dla takich wywołań, z których można
utworzyć poprawne wywołanie \texttt{output}. Sama wariadyczność
nie oznacza, że każda liczba i każdy rodzaj argumentów będą
akceptowane przez funkcję docelową.

\subsection{Przekazywanie przez wartość i przez referencję}

\begin{lstlisting}
template<class... Ts>
void by_value(Ts... xs);            // may copy/move

template<class... Ts>
void by_const_ref(const Ts&... xs); // read-only; no copies

template<class... Ts>
void by_ref(Ts&... xs);             // needs lvalues
\end{lstlisting}

Wybór postaci parametrów ma te same skutki co w zwykłych funkcjach.
W szczególności \texttt{Ts... xs} tworzy osobne parametry przez wartość,
a \texttt{const Ts\&... xs} zwykle pozwala uniknąć niepotrzebnego
kopiowania. Przypadek \texttt{Ts\&\&... xs} omówimy dokładniej
w sekcji o \emph{perfect forwarding}.

\section{Rekurencyjne przetwarzanie pakietu: metoda C++11}
\label{sec:variadic-rekurencja}

Przed C++17 często rozbijano pakiet na \emph{pierwszy argument}
i \emph{pozostałe argumenty}, a następnie wywoływano funkcję
rekurencyjnie. Potrzebny był osobny przypadek kończący rekurencję.

\begin{lstlisting}
#include <iostream>

void print_recursive() {
    std::cout << '\n';        // base case
}

template<class First, class... Rest>
void print_recursive(const First& first,
                     const Rest&... rest) {
    std::cout << first << ' ';
    print_recursive(rest...); // one argument fewer
}

int main() {
    print_recursive(1, "two", 3.0);
}
\end{lstlisting}

Przy trzech argumentach instancjonowane są kolejno wersje
z pakietami o długościach $3$, $2$, $1$ oraz przypadek końcowy.
Jest to \emph{rekurencja szablonowa}; nie należy automatycznie
utożsamiać jej z dynamiczną rekurencją wykonywaną w programie.
Kompilator może wygenerować zwykłą sekwencję operacji.

\paragraph{Kiedy stosować?}
Rekurencyjne rozkładanie pakietu nadal bywa przydatne, jeśli kolejne
kroki mają odrębne warunki lub należy budować typy rekurencyjne.
Jeżeli wykonujemy tę samą operację dla wszystkich argumentów,
w C++17 i nowszych zazwyczaj lepszym wyborem jest \emph{fold expression}.

\section{Fold expressions (C++17)}

\begin{center}\small
\begin{tabularx}{\textwidth}{@{}p{0.35\textwidth}X@{}}
\toprule
Forma złożenia & Sposób rozwinięcia \\
\midrule
\code{(... + args)} & Złożenie lewostronne: kolejne dodawania łączone od lewej. \\
\code{(args + ...)} & Złożenie prawostronne: kolejne dodawania łączone od prawej. \\
\code{(init + ... + args)} & Złożenie z wartością początkową; działa także dla pustego pakietu. \\
\code{(f(args), ...)} & Wywołuje \code{f} dla każdego argumentu od lewej do prawej dzięki operatorowi przecinka. \\
\texttt{(args \&\& ...)} & Łączy warunki z krótkim spięciem; dla pustego pakietu daje \code{true}. \\
\bottomrule
\end{tabularx}
\end{center}
\label{sec:variadic-fold}

\emph{Fold expression} redukuje pakiet wyrażeń z użyciem jednego
operatora binarnego, np. \texttt{+}, \texttt{\&\&}, \texttt{||}
lub operatora przecinka. Nie wymaga jawnego przypadku
rekurencyjnego.

\subsection{Cztery postacie składni}

Niech \texttt{xs} oznacza pakiet, $\circ$ ustalony operator,
a $z$ wartość początkową. Istnieją cztery postacie składni:

\begin{center}
\small
\begin{tabular}{|p{0.39\textwidth}|p{0.47\textwidth}|}
\hline
\textbf{Postać C++} & \textbf{Sposób łączenia dla trzech argumentów} \\
\hline
\texttt{(xs op ...)} & $x_1\circ(x_2\circ x_3)$ (prawostronne, bez inicjalizatora) \\
\hline
\texttt{(... op xs)} & $(x_1\circ x_2)\circ x_3$ (lewostronne, bez inicjalizatora) \\
\hline
\texttt{(xs op ... op z)} & $x_1\circ(x_2\circ(x_3\circ z))$ (prawostronne, z inicjalizatorem) \\
\hline
\texttt{(z op ... op xs)} & $((z\circ x_1)\circ x_2)\circ x_3$ (lewostronne, z inicjalizatorem) \\
\hline
\end{tabular}
\end{center}

W tabeli \texttt{op} jest jedynie oznaczeniem miejsca na rzeczywisty
operator, np. \texttt{+}; nie jest słowem kluczowym języka.
\emph{Nawiasy całego fold expression są obowiązkowe}.
Dwa operatory występujące w fold z inicjalizatorem muszą być takie same.

\subsection{Suma, iloczyn i warunki logiczne}

\begin{lstlisting}
#include <concepts>

template<class... Ts>
    requires (std::integral<Ts> && ...)
constexpr auto sum(Ts... xs) {
    return (0 + ... + xs);       // left fold with identity 0
}

template<class... Ts>
constexpr bool all_true(Ts... xs) {
    return (true && ... && static_cast<bool>(xs));
}

template<class... Ts>
constexpr bool any_true(Ts... xs) {
    return (false || ... || static_cast<bool>(xs));
}

static_assert(sum() == 0);
static_assert(sum(2, 3, 4) == 9);
static_assert(all_true());
static_assert(!any_true());
static_assert(all_true(1, 2, 3));
\end{lstlisting}

Wyrażenie \texttt{(0 + ... + xs)} może być używane także dla
pustego pakietu. Początkowe \texttt{0} wpływa jednak na typ wyniku
i ewentualne promocje arytmetyczne; nie jest to uniwersalna
implementacja sumy dowolnych typów liczbowych.

\subsection{Lewostronne i prawostronne składanie nie zawsze są równe}

Dla operatora odejmowania kolejność nawiasowania zmienia wynik:

\begin{lstlisting}
template<class... Ts>
constexpr auto left_sub(Ts... xs) {
    return (100 - ... - xs);
}

template<class... Ts>
constexpr auto right_sub(Ts... xs) {
    return (xs - ... - 100);
}

static_assert(left_sub(10, 3) == 87);  // (100 - 10) - 3
static_assert(right_sub(10, 3) == 107); // 10 - (3 - 100)
\end{lstlisting}

Fold określa \emph{strukturę nawiasów}, ale nie należy mylić tej
własności z ogólną gwarancją kolejności wykonywania wszystkich
podwyrażeń dla dowolnego operatora. Reguły ewaluacji wynikają
również z semantyki danego operatora.

\subsection{Pusty pakiet i element neutralny}

Dla \emph{jednoargumentowego} fold z pustym pakietem C++ dopuszcza
zaledwie trzy operatory: \texttt{\&\&} zwraca \texttt{true},
\texttt{||} zwraca \texttt{false}, a \texttt{,} daje
wyrażenie typu \texttt{void}. Pusty pakiet w
\texttt{(... + xs)} nie jest poprawny. Stosując wersję
z wartością początkową, możemy jawnie obsłużyć pusty pakiet:

\begin{lstlisting}
template<class... Ts>
constexpr auto product(Ts... xs) {
    return (1 * ... * xs);
}

static_assert(product() == 1);
static_assert(product(2, 3, 4) == 24);
\end{lstlisting}

\section{Operator przecinka i kolejność wykonywania}
\label{sec:variadic-przecinek}

Operator przecinka jest wygodny, gdy dla każdego elementu pakietu
chcemy wykonać instrukcję o efektach ubocznych. W odróżnieniu
od kolejności ewaluacji zwykłych argumentów funkcji, jego
wbudowana postać gwarantuje wykonanie lewej strony przed prawą.

\begin{lstlisting}
#include <iostream>

template<class... Ts>
void print_all(const Ts&... xs) {
    ((std::cout << xs << ' '), ...);
    std::cout << '\n';
}

int main() {
    print_all(10, 2.5, "hello");
    print_all();                // only a newline
}
\end{lstlisting}

Zewnętrzne nawiasy tworzą fold, a wewnętrzne oddzielają
wyrażenie wypisujące od operatora \texttt{,}.

\paragraph{Kolejność argumentów zwykłej funkcji.}

\begin{lstlisting}
template<class... Ts>
void sink(Ts...);

// sink(change_a(), change_b(), change_c());
\end{lstlisting}

Samo rozwinięcie \texttt{...} w liście argumentów funkcji
\emph{nie} gwarantuje, że \texttt{change\_a()} zostanie obliczone
przed \texttt{change\_b()}. W C++17 i nowszych wywołania
te nie przeplatają swoich efektów, ale ich wzajemna kolejność
nie jest określona. Gdy kolejność jest istotna, można użyć
operatora przecinka lub innej konstrukcji dającej odpowiednie
gwarancje sekwencjonowania.

\section{Pakiety typów i metaprogramowanie}
\label{sec:variadic-type-list}

Wariadyczność może dotyczyć nie tylko funkcji, lecz także
\emph{szablonów klas}. W takim przypadku pakiet przechowuje
listę typów, a instancjonowanie klasy odbywa się podczas kompilacji.

\subsection{Własna lista typów}

\begin{lstlisting}
#include <cstddef>
#include <type_traits>

template<class... Ts>
struct type_list {};

template<class List>
struct list_size;

template<class... Ts>
struct list_size<type_list<Ts...>>
    : std::integral_constant<std::size_t, sizeof...(Ts)> {};

using Types = type_list<int, double, char>;
static_assert(list_size<Types>::value == 3);
static_assert(list_size<type_list<>>::value == 0);
\end{lstlisting}

Jest to przykład \emph{częściowej specjalizacji} szablonu klasy.
Wzorzec \texttt{type\_list<Ts...>} dopasowuje dowolną listę
typów, a \texttt{sizeof...(Ts)} oblicza jej długość.

\subsection{Wydobywanie pierwszego typu}

\begin{lstlisting}
template<class List>
struct first_type;   // intentionally undefined for empty list

template<class T, class... Rest>
struct first_type<type_list<T, Rest...>> {
    using type = T;
};

using X = first_type<type_list<long, float>>::type;
static_assert(std::is_same_v<X, long>);
\end{lstlisting}

Rozkład na \texttt{T} oraz \texttt{Rest...} jest odpowiednikiem
rekurencyjnego rozdzielenia argumentów funkcji na pierwszy
i pozostałe. Dla pustej listy \texttt{type\_list<>} nazwa
\texttt{first\_type<...>::type} celowo nie istnieje.

\subsection{Czy wszystkie typy spełniają własność?}

\begin{lstlisting}
#include <type_traits>

template<class... Ts>
inline constexpr bool all_integral_v =
    (std::is_integral_v<Ts> && ...);

template<class... Ts>
inline constexpr std::size_t count_integral_v =
    (std::size_t{0} + ... +
        (std::is_integral_v<Ts> ? 1u : 0u));

static_assert(all_integral_v<int, long, char>);
static_assert(!all_integral_v<int, double>);
static_assert(count_integral_v<int, double, char> == 2);
\end{lstlisting}

Jest to częsty wzorzec: pakiet typów łączymy z cechami
typów (\emph{type traits}) i operatorami logicznymi fold.

\section{Pakiety wartości i parametry nietypowe}

Pakiet może także zawierać wartości znane w czasie kompilacji,
a nie typy. W szablonie klasy służą do tego
\emph{nietypowe parametry szablonu} (\emph{non-type template parameters}).

\begin{lstlisting}
#include <array>
#include <cstddef>

template<int... Ns>
struct int_pack {
    static constexpr std::size_t size = sizeof...(Ns);
    static constexpr int sum = (0 + ... + Ns);
};

static_assert(int_pack<>::sum == 0);
static_assert(int_pack<2, 4, 6>::sum == 12);
\end{lstlisting}

Od C++17 można zadeklarować pakiet \texttt{auto...}:

\begin{lstlisting}
template<auto... Values>
struct value_pack {
    static constexpr std::size_t size = sizeof...(Values);
};

static_assert(value_pack<42, 'x', true>::size == 3);
\end{lstlisting}

Dla \texttt{auto...} dopuszczalność konkretnej wartości jako
argumentu szablonu nadal podlega regułom parametrów nietypowych;
nie można przekazać do niego dowolnej wartości z czasu wykonania.

\section{Przekazywanie doskonałe (\textit{perfect forwarding})}
\label{sec:variadic-forwarding}

Jednym z najważniejszych zastosowań pakietów jest budowanie
\emph{funkcji pośredniczących}, które przekazują argumenty do
innej funkcji lub konstruktora, możliwie wiernie zachowując
ich kategorię wartości (\emph{value category}).

\subsection{Różnica między lvalue a rvalue}

\begin{lstlisting}
#include <iostream>

void use(int&)  { std::cout << "lvalue\n"; }
void use(int&&) { std::cout << "rvalue\n"; }

int main() {
    int x = 10;
    use(x);       // lvalue
    use(20);      // rvalue
}
\end{lstlisting}

Parametr funkcji, choćby zadeklarowany jako \texttt{T\&\&}, jest
\emph{nazwanym obiektem}; użycie jego nazwy w wyrażeniu ma
kategorię \emph{lvalue}. Dlatego samo \texttt{use(arg)}
nie zachowa kategorii pierwotnego argumentu.

\subsection{Referencje przekazujące i \texttt{std::forward}}

\begin{lstlisting}
#include <utility>

template<class T>
void wrapper(T&& arg) {
    use(std::forward<T>(arg));
}

void forwarding_example() {
    int y = 0;
    wrapper(y);  // T = int&, forwarded as lvalue
    wrapper(30); // T = int, forwarded as rvalue
}
\end{lstlisting}

\texttt{T\&\&} jest \emph{referencją przekazującą}
(ang. \emph{forwarding reference}), gdy \texttt{T} jest
odpowiednio dedukowanym, niekwalifikowanym parametrem typu.
Działa tu \emph{składanie referencji}:

\[
T\&\;\& \longrightarrow T\&,
\qquad
T\&\;\&\& \longrightarrow T\&,
\qquad
T\&\&\;\&\& \longrightarrow T\&\&.
\]

W zapisie \texttt{std::forward<T>(arg)} jawny argument
\texttt{T} jest konieczny, by odtworzyć pierwotną kategorię.
\texttt{std::forward} nie wykonuje bezwarunkowego przeniesienia:
przepuszcza lvalue jako lvalue, a rvalue jako rvalue.

\subsection{Pakiet referencji przekazujących}

\begin{lstlisting}
#include <functional>
#include <utility>

template<class F, class... Args>
decltype(auto) call(F&& f, Args&&... args) {
    return std::invoke(
        std::forward<F>(f),
        std::forward<Args>(args)...
    );
}
\end{lstlisting}

Tutaj rozwijany jest \emph{cały wzorzec}
\texttt{std::forward<Args>(args)}, a nie tylko nazwa
\texttt{args}. Każdy element używa właściwego typu z pakietu
\texttt{Args}.
\texttt{std::invoke} potrafi wywoływać zwykłe funkcje,
obiekty funkcyjne i odpowiednio obsłużone wskaźniki do składowych.
\texttt{decltype(auto)} zachowuje typ zwracany przez docelowe
wywołanie, w tym referencję, jeżeli docelowa funkcja ją zwraca.

\begin{lstlisting}
#include <string>

std::string join(std::string a, std::string b) {
    return a + b;
}

auto result = call(join, std::string{"ab"}, "cd");
// result == "abcd"
\end{lstlisting}

\subsection{Czego nie robić: \texttt{std::move} i nazwy parametrów}

\begin{lstlisting}
template<class T>
void wrong_1(T&& arg) {
    use(arg);              // arg is always an lvalue here
}

template<class T>
void wrong_2(T&& arg) {
    use(std::move(arg));   // forces an rvalue even for lvalues
}
\end{lstlisting}

\texttt{std::move} to rzutowanie umożliwiające potraktowanie
obiektu jak rvalue, a nie polecenie ``zachowaj kategorię''.
Stosuj je, gdy rzeczywiście przekazujesz obiekt do operacji,
która może pobrać z niego zasób. W uniwersalnym wrapperze
używaj \texttt{std::forward}.

\paragraph{Ograniczenia żywotności.}
Referencje przekazujące \emph{nie przedłużają automatycznie}
żywotności obiektu do czasu dowolnego późniejszego wykonania.
Przechowanie referencji do argumentu tymczasowego w strukturze,
która przeżyje wywołanie, może prowadzić do wskaźników lub
referencji wiszących. Ten sam problem dotyczy między innymi
\texttt{std::forward\_as\_tuple}.

\section{Tworzenie obiektów z dowolnych argumentów}
\label{sec:variadic-emplace}

\subsection{Własna fabryka obiektów}

\begin{lstlisting}
#include <memory>
#include <utility>

template<class T, class... Args>
std::unique_ptr<T> create(Args&&... args) {
    return std::make_unique<T>(
        std::forward<Args>(args)...
    );
}

struct Point {
    int x, y;
    Point(int a, int b) : x(a), y(b) {}
};

auto point = create<Point>(4, 7);
\end{lstlisting}

\texttt{std::make\_unique} już dostarcza taką funkcjonalność,
więc \texttt{create} ma tu charakter dydaktyczny. Pokazuje
jednak wzorzec używany w fabrykach oraz funkcjach
\texttt{emplace} i \texttt{make}.

\subsection{Tworzenie kontenera bez wymogu kopiowalności}

\begin{lstlisting}
#include <concepts>
#include <string>
#include <utility>
#include <vector>

template<class T, class... Args>
    requires (std::constructible_from<T, Args&&> && ...)
std::vector<T> make_vector(Args&&... args) {
    std::vector<T> result;
    result.reserve(sizeof...(Args));
    (result.emplace_back(std::forward<Args>(args)), ...);
    return result;
}

auto words = make_vector<std::string>("A", "BB", "CCC");
\end{lstlisting}

Operator przecinka daje tu jednoznaczną kolejność wywołań
\texttt{emplace\_back}. Funkcja przyjmuje także pusty pakiet
argumentów. W przeciwieństwie do budowania wektora z listy
inicjalizacyjnej pozwala łatwo przekazywać do kolejnych
\texttt{emplace\_back} argumenty przenoszalne, ale niekopiowalne.
Dla elementów konstruowanych z więcej niż jednego parametru
wymagałaby innego interfejsu grupowania argumentów.

\section{\texttt{std::tuple}, \texttt{std::apply} i \texttt{std::index\_sequence}}
\label{sec:variadic-tuple}

\subsection{Heterogeniczna krotka \texttt{std::tuple}}

\texttt{std::tuple<Ts...>} to kontener o strukturze opisanej
listą typów znaną w czasie kompilacji. Każda pozycja może
mieć inny typ:

\begin{lstlisting}
#include <tuple>
#include <string>

std::tuple<int, std::string, double> t{7, "Ada", 3.5};

int n = std::get<0>(t);
const std::string& name = std::get<1>(t);
\end{lstlisting}

W przeciwieństwie do \texttt{std::vector<T>}, krotka nie
przechowuje dowolnej liczby elementów jednego typu, lecz
\emph{ustaloną na etapie kompilacji} listę typów.

\subsection{Rozwinięcie krotki do argumentów funkcji}

Funkcja \texttt{std::apply} (C++17) wywołuje funkcję,
przekazując elementy obiektu \emph{tuple-like} jako osobne
argumenty:

\begin{lstlisting}
#include <iostream>
#include <tuple>

void apply_example() {
    std::apply([](const auto&... xs) {
        ((std::cout << xs << ' '), ...);
        std::cout << '\n';
    }, t);
}
\end{lstlisting}

\texttt{std::apply} oszczędza ręcznego konstruowania
sekwencji indeksów, ale warto wiedzieć, jak taki mechanizm
można zaimplementować.

\subsection{Pakiety indeksów: \texttt{std::index\_sequence}}

\begin{lstlisting}
#include <cstddef>
#include <iostream>
#include <tuple>
#include <utility>

template<class Tuple, std::size_t... I>
void print_tuple_impl(const Tuple& t,
                      std::index_sequence<I...>) {
    ((std::cout << std::get<I>(t) << ' '), ...);
    std::cout << '\n';
}

template<class... Ts>
void print_tuple(const std::tuple<Ts...>& t) {
    print_tuple_impl(t, std::index_sequence_for<Ts...>{});
}

void tuple_example() {
    print_tuple(std::make_tuple(1, "abc", 3.0));
}
\end{lstlisting}

\texttt{std::index\_sequence<0,1,2>} jest specjalizacją
szablonu przyjmującego pakiet indeksów całkowitych.
\texttt{std::index\_sequence\_for<Ts...>} generuje indeksy
$0,1,\ldots,n-1$ dla pakietu długości $n$.
Wzorzec \texttt{std::get<I>(t)...} rozwija wyrażenie
osobno dla każdego indeksu.

\subsection{Równoległe rozwijanie pakietów}

W jednym wzorcu można użyć kilku pakietów. Muszą mieć
\emph{tę samą długość}, jeśli są rozwijane równolegle:

\begin{lstlisting}
#include <tuple>
#include <type_traits>
#include <utility>

template<class A, class B>
struct zipped_types;

template<class... Ts, class... Us>
struct zipped_types<std::tuple<Ts...>, std::tuple<Us...>> {
    static_assert(sizeof...(Ts) == sizeof...(Us));
    using type = std::tuple<std::pair<Ts, Us>...>;
};

using Z = zipped_types<
    std::tuple<int, char>,
    std::tuple<double, bool>
>::type;

static_assert(std::is_same_v<
    Z, std::tuple<std::pair<int, double>,
                  std::pair<char, bool>>>);
\end{lstlisting}

Elementy \texttt{Ts} i \texttt{Us} są łączone pozycjami:
$\bigl(T_1,U_1\bigr),\ldots,\bigl(T_n,U_n\bigr)$.
Nie jest to iloczyn kartezjański pakietów.

\section{Dziedziczenie wariadyczne i \texttt{std::visit}}
\label{sec:variadic-overloaded}

Pakiet \texttt{Bases...} można rozwinąć na liście baz.
Z kolei zapis \texttt{using Bases::operator()...}
wprowadza ich operatory wywołania do jednej klasy.
W ten sposób otrzymujemy zestaw przeciążeń.

\begin{lstlisting}
#include <iostream>
#include <string>
#include <variant>

template<class... Fs>
struct overloaded : Fs... {
    using Fs::operator()...;
};

template<class... Fs>
overloaded(Fs...) -> overloaded<Fs...>; // C++17 deduction guide

int main() {
    std::variant<int, std::string> value = "hello";

    std::visit(overloaded{
        [](int n) { std::cout << "int: " << n << '\n'; },
        [](const std::string& s) {
            std::cout << "string: " << s << '\n';
        }
    }, value);
}
\end{lstlisting}

Jest to standardowy idiom C++17: każda lambda definiuje
osobny \texttt{operator()}, a dziedziczenie wariadyczne
łączy je w jeden obiekt wywoływalny. Gdy stosujemy ten wzorzec
w wielu miejscach, warto umieścić \texttt{overloaded} w
niewielkim nagłówku narzędziowym.

\section{Ograniczanie pakietów za pomocą \texttt{concepts} (C++20)}
\label{sec:variadic-concepts}

Nie każda operacja ma sens dla każdej możliwej kombinacji
typów. W C++20 można czytelnie określać wymagania dla
\emph{wszystkich} elementów pakietu.

\begin{lstlisting}
#include <concepts>
#include <type_traits>

template<class... Ts>
    requires (std::integral<Ts> && ...)
constexpr bool all_nonnegative(Ts... xs) {
    return ((xs >= 0) && ...);
}

static_assert(all_nonnegative(0, 2, 7));
static_assert(!all_nonnegative(0, -1, 7));

// all_nonnegative(3.14); // rejected by constraints
\end{lstlisting}

Możemy również ograniczać pakiet w odniesieniu do
konkretnego typu docelowego:

\begin{lstlisting}
#include <concepts>

template<class T, class... Args>
concept ConstructibleWith =
    std::constructible_from<T, Args...>;

static_assert(ConstructibleWith<std::string, const char*>);
static_assert(!ConstructibleWith<int, std::string>);
\end{lstlisting}

Pierwszy przykład stosuje fold nad \emph{warunkami dla pojedynczych
typów}; drugi sprawdza \emph{jedną własność całej sekwencji
typów}, tj. czy odpowiedni konstruktor przyjmuje argumenty
\texttt{Args...} w tej kolejności. Tych dwóch mechanizmów
nie należy utożsamiać.

\subsection{\texttt{if constexpr} i struktura pakietu}

\begin{lstlisting}
#include <iostream>

template<class... Ts>
void report(const Ts&... xs) {
    if constexpr (sizeof...(Ts) == 0) {
        std::cout << "empty\n";
    } else {
        ((std::cout << xs << ' '), ...);
        std::cout << '\n';
    }
}
\end{lstlisting}

\texttt{if constexpr} pozwala w zależności od długości pakietu
wybierać gałąź w czasie kompilacji. W dobrze sformułowanym
szablonie odrzucona gałąź nie jest instancjonowana dla danej
specjalizacji. To użyteczne przy przypadkach brzegowych,
które nie mają sensu dla pustego pakietu.

\section{Ograniczenia, pułapki i diagnostyka}
\label{sec:variadic-pulapki}

\subsection{Pakiet typów nie jest pakietem argumentów funkcji}

\begin{lstlisting}
template<class... Ts>
void f(Ts... args) {
    // sizeof...(Ts): number of types
    // sizeof...(args): number of function parameters
    // args[0];  // ERROR: args is not an array
}
\end{lstlisting}

Do wybierania argumentu na ustalonej pozycji można użyć
\texttt{std::tuple} i \texttt{std::get} albo wywoływać funkcję
przez \texttt{std::apply}. Sam pakiet \texttt{args...}
nie jest indeksowanym kontenerem.

\subsection{Brak zgodności typów po rozwinięciu}

\begin{lstlisting}
template<class... Ts>
auto add(Ts... xs) {
    return (0 + ... + xs);
}

// add(1, std::string{"x"}); // ERROR: unsuitable operator+
\end{lstlisting}

Pakiet nie omija systemu typów. Każde wyrażenie powstałe po
jego rozwinięciu musi być poprawne, podobnie jak zwykły kod
C++. Jeżeli szablon ma określone wymagania, najlepiej
zapisać je za pomocą \texttt{requires}.

\subsection{Niewłaściwe przekazywanie kategorii wartości}

\begin{lstlisting}
template<class F, class... Args>
void bad_call(F&& f, Args&&... args) {
    std::forward<F>(f)(args...); // loses rvalue categories
}

// Correct:
// std::forward<F>(f)(std::forward<Args>(args)...);
\end{lstlisting}

Dla typów niekopiowalnych taki błąd często ujawnia się
natychmiast podczas kompilacji. Dla typów kopiowalnych
może spowodować jedynie dodatkowe kopie albo zmianę
wybranego przeciążenia.

\subsection{Wnioskowanie typów i \texttt{\{...\}}}

\begin{lstlisting}
#include <initializer_list>

template<class... Args>
void foo(Args&&...) {}

void braced_example() {
    // foo({1, 2, 3}); // ERROR: cannot deduce Args
    foo(std::initializer_list<int>{1, 2, 3}); // OK
}
\end{lstlisting}

Sama lista ujęta w nawiasy klamrowe \texttt{\{...\}} nie ma
typu, który dowolny pakiet \texttt{Args\&\&...} potrafiłby
wydedukować. Potrzebny jest określony typ parametru,
np. \texttt{std::initializer\_list}\allowbreak\texttt{<int>}, lub jawnie
utworzony obiekt.

\subsection{Ograniczenia położenia pakietu}

W \emph{głównej definicji szablonu klasy} pakiet parametrów
szablonu musi być ostatni. Dla szablonów funkcji istnieją
bardziej elastyczne reguły, ale kolejne parametry muszą
być dedukowalne lub mieć argumenty domyślne, stosownie do
kontekstu:

\begin{lstlisting}
// ERROR for a primary class template:
// template<class... Ts, class U> struct invalid;

template<class U, class... Ts>
struct valid {};
\end{lstlisting}

\subsection{Skomplikowany kod zamiast prostego rozwiązania}

Szablon wariadyczny to narzędzie, a nie obowiązkowy styl.
Jeżeli argumenty mają \emph{jednolity typ i ich liczba
zmienia się w czasie wykonania}, lepszy bywa
\texttt{std::span<T>}, \texttt{std::vector<T>} lub para
iteratorów. Pakiet ma liczbę elementów i strukturę typów
ustalaną na etapie kompilacji; nie zastępuje kontenera.

\section{Variadic templates a funkcje wariadyczne w stylu C}

Historyczne funkcje z deklaracją \texttt{void f(int, ...)}
korzystają z mechanizmu argumentów wariadycznych języka C,
a nie z pakietu parametrów szablonu C++:

\begin{lstlisting}
#include <cstdio>

int main() {
    std::printf("%d %f\n", 7, 2.5); // C-style variadic call
}
\end{lstlisting}

W formacie C zgodność argumentów ze specyfikatorami
formatowania jest obowiązkiem programisty (niektóre
kompilatory potrafią wykrywać część pomyłek). W szablonach
wariadycznych typy argumentów pozostają dostępne
kompilatorowi, dzięki czemu można wymuszać ograniczenia
i stosować przeciążanie oraz \texttt{std::forward}.
Nie należy jednak wyciągać z tego wniosku, że dowolny
szablon automatycznie bezbłędnie formatuje argumenty.

\section{Pełny przykład: prosty rejestrator zdarzeń}
\label{sec:variadic-projekt}

Zastosujmy omówione mechanizmy do funkcji tworzącej tekstowy
komunikat z dowolnej liczby argumentów możliwych do
wypisania do strumienia. Zamiast ręcznie doklejać spacje
w wielu przeciążeniach, użyjemy fold nad operatorem przecinka.

\begin{lstlisting}
#include <iostream>
#include <sstream>
#include <string>

// Each argument must support insertion into std::ostream.
template<class... Args>
std::string make_message(const Args&... args) {
    std::ostringstream out;
    ((out << args), ...);
    return out.str();
}

template<class... Args>
void log_line(const Args&... args) {
    std::cout << "[LOG] " << make_message(args...) << '\n';
}

int main() {
    log_line("count=", 3, ", ratio=", 0.75);
    log_line("finished");
    log_line();
}
\end{lstlisting}

\paragraph{Analiza.}
Dla argumentów \texttt{("count=",3,", ratio=",0.75)}
pakiet zawiera cztery elementy. Wyrażenie
\texttt{((out << args), ...)} wstawia je kolejno do strumienia,
a \texttt{make\_message} zwraca powstały napis. Dla
pustego pakietu fold z operatorem przecinka jest poprawny,
a wynikowy napis jest pusty. W wersji produkcyjnej należałoby
rozważyć wymagania dotyczące formatowania, zachowanie w
obliczu wyjątków, współbieżność oraz ewentualne skorzystanie
z narzędzia formatującego wspieranego przez używany kompilator.

\section{Zestawienie: kiedy używać poszczególnych mechanizmów?}

\begin{center}
\small
\begin{tabular}{|p{0.40\textwidth}|p{0.46\textwidth}|}
\hline
\textbf{Cel} & \textbf{Typowe rozwiązanie} \\
\hline
Dowolna liczba argumentów o różnych typach & \texttt{template<class... Ts>} i \texttt{Ts...} \\
\hline
Wykonanie tej samej operacji dla każdego argumentu & Fold z operatorem \texttt{,} (C++17) \\
\hline
Suma, iloczyn lub wspólny warunek & Fold z \texttt{+}, \texttt{*}, \texttt{\&\&} lub \texttt{||} \\
\hline
Przekazanie argumentów dalej z zachowaniem kategorii & \texttt{Args\&\&...}, \texttt{std::forward<Args>(args)...} \\
\hline
Zapis heterogenicznych wartości & \texttt{std::tuple<Ts...>} \\
\hline
Wywołanie funkcji elementami krotki & \texttt{std::apply} (C++17) \\
\hline
Rozwinięcie krotki przy pomocy indeksów & \texttt{std::index\_sequence} \\
\hline
Lista typów do obliczeń podczas kompilacji & Wariadyczny szablon klasy i specjalizacje \\
\hline
Ograniczenie do odpowiednich typów & Fold z koncepcjami i \texttt{requires} (C++20) \\
\hline
Zmienna długość listy jednolitych danych w czasie wykonania & \texttt{std::vector} lub \texttt{std::span}, nie pakiet \\
\hline
\end{tabular}
\end{center}

\section{Podsumowanie i zadania do samodzielnego sprawdzenia}

\begin{enumerate}
    \item \texttt{class... Ts} deklaruje \emph{pakiet typów},
    a \texttt{Ts... xs} wprowadza \emph{pakiet parametrów funkcji}.
    \item \texttt{sizeof...(Ts)} oblicza liczbę elementów pakietu
    podczas kompilacji.
    \item \texttt{pattern...} rozwija cały \emph{wzorzec},
    nie tylko pojedynczy identyfikator.
    \item Fold expressions (od C++17) zastępują wiele prostych
    rekurencyjnych implementacji z C++11.
    \item Dla pustego pakietu istotne są element neutralny i
    reguły operatora; nie każda forma fold jest poprawna.
    \item W wrapperach zachowuj kategorię wartości argumentów,
    używając referencji przekazujących i \texttt{std::forward}.
    \item \texttt{std::tuple}, \texttt{std::apply} i
    \texttt{std::index\_sequence} współpracują z pakietami,
    umożliwiając operacje na heterogenicznych danych.
    \item W C++20 używaj \texttt{concepts}/\texttt{requires},
    aby błędne wywołania odrzucać na granicy interfejsu.
\end{enumerate}

\paragraph{Pytania kontrolne.}
Spróbuj przewidzieć wyniki poniższych wyrażeń bez uruchamiania
programu. Następnie zweryfikuj je za pomocą \texttt{static\_assert}:

\begin{lstlisting}
// static_assert(sizeof...(int) == 1); // ERROR: int is not a pack

// For the functions defined above:
static_assert(product() == 1);
static_assert(left_sub(10, 3) == 87);
static_assert(right_sub(10, 3) == 107);
static_assert(count_integral_v<int, double, char> == 2);
\end{lstlisting}

Pierwsza linia jest \emph{celowo niepoprawna}: operator
\texttt{sizeof...} wymaga nazwy pakietu, a nie zwykłego typu.

% Wykorzystane punkty odniesienia (bez kopiowania tresci):
% https://en.cppreference.com/w/cpp/language/parameter_pack
% https://en.cppreference.com/w/cpp/language/fold
% https://en.cppreference.com/w/cpp/utility/forward
% https://en.cppreference.com/w/cpp/utility/apply
% https://en.cppreference.com/w/cpp/utility/integer_sequence


\chapter[Variadic templates w C++]{Szablony o zmiennej liczbie argumentów (\textit{variadic templates})}
\label{chap:variadic-templates-cpp}

\section{Idea i zastosowania}

Zwykły szablon funkcji przyjmuje z góry ustaloną liczbę parametrów
szablonu. \emph{Variadic template} (szablon o zmiennej liczbie parametrów)
może natomiast przyjmować \emph{dowolną liczbę} argumentów szablonu,
również zero. Mechanizmem, który to umożliwia, jest
\emph{pakiet parametrów} (ang. \emph{parameter pack}).

\begin{lstlisting}
template<class T>
void single(const T&) {}        // one type

template<class... Ts>
void many(const Ts&...) {}      // zero or more types

void demo() {
    many();                     // zero arguments
    many(10);                   // one argument
    many(10, 3.5, "abc");      // three arguments
}
\end{lstlisting}

Pakiety parametrów są ważną częścią nowoczesnego C++.
Wykorzystują je m.in. klasy \texttt{std::tuple} i \texttt{std::variant}
oraz fabryki \texttt{std::make\_unique} i \texttt{std::make\_shared}.
Nie trzeba dzięki nim pisać przeciążeń dla każdej
liczby argumentów.

\paragraph{Wymagane wersje języka.}
Pakiety parametrów i perfect forwarding pojawiły się w C++11.
\emph{Fold expressions} dodano w C++17, a koncepcje
(\texttt{concepts}) i klauzulę \texttt{requires} w C++20.
Dalej zakładamy C++20.

\section{Trzy kropki: deklaracja pakietu a rozwinięcie}

\begin{center}\small
\begin{tabularx}{\textwidth}{@{}p{0.35\textwidth}X@{}}
\toprule
Zapis & Działanie \\
\midrule
\code{typename... Ts} & Deklaruje pakiet parametrów typowych szablonu. \\
\code{Ts...} & Rozwija pakiet typów w miejscu, które dopuszcza taki wzorzec. \\
\code{args...} & Rozwija pakiet argumentów, tworząc listę wyrażeń. \\
\code{sizeof...(Ts)} & Podaje liczbę elementów pakietu w czasie kompilacji. \\
\code{std::forward<Ts>(args)...} & Przekazuje każdy argument z zachowaniem jego kategorii wartości. \\
\bottomrule
\end{tabularx}
\end{center}
\label{sec:variadic-trzy-kropki}

W różnych miejscach token \texttt{...} oznacza różne operacje.
Rozróżnienie to stanowi podstawę poprawnego używania szablonów
wariadycznych.

\begin{lstlisting}
template<class... Ts>             // declare type pack Ts
void f(Ts... xs) {                // declare function parameter pack xs
    // ...
}
\end{lstlisting}

Deklaracja \texttt{class... Ts} wprowadza pakiet \emph{typów}.
Zapis \texttt{Ts... xs} rozwija pakiet typów i wprowadza odpowiadający
mu pakiet parametrów funkcji. Wywołanie
\texttt{f(1, 2.0, 'x')} tworzy w uproszczeniu funkcję o parametrach
\texttt{(int, double, char)}.

\subsection{Rozwinięcie pakietu (\textit{pack expansion})}

Rozwinięcie składa się ze \emph{wzorca} (\emph{pattern}) oraz następującego
po nim \texttt{...}. Dla każdego elementu pakietu kompilator podstawia
kolejny element do całego wzorca:

\begin{lstlisting}
template<class... Ts>
void demo(Ts... xs) {
    g(xs...);                // g(x1, x2, ..., xn)
    h((xs * 2)...);          // h(x1*2, x2*2, ..., xn*2)
}
\end{lstlisting}

W drugim wywołaniu wzorcem jest \texttt{(xs * 2)}, a \emph{nie}
pojedyncza zmienna \texttt{xs}. Ważne: \texttt{xs...} nie jest
kontenerem ani tablicą; nie można po nim bezpośrednio iterować zwykłą
pętlą \texttt{for}. To składnia kompilatora służąca do generowania
listy wyrażeń.

\paragraph{Pakiet typów a pakiet wartości.}
W przykładzie \texttt{Ts} reprezentuje typy, natomiast
\texttt{xs} to konkretne parametry funkcji. Jeden typ może wystąpić
w pakiecie wielokrotnie, a każdy element pakietu wartości zachowuje
swój właściwy typ.

\subsection{Liczba elementów: \texttt{sizeof...}}

Operator \texttt{sizeof...} daje liczbę elementów pakietu jako
stałą czasu kompilacji. Nie należy go mylić z operatorem
\texttt{sizeof} mierzącym rozmiar typu lub obiektu.

\begin{lstlisting}
#include <cstddef>

template<class... Ts>
constexpr std::size_t count(Ts&&... xs) {
    static_assert(sizeof...(Ts) == sizeof...(xs));
    return sizeof...(xs);
}

static_assert(count() == 0);
static_assert(count(2, 3.5, 'a') == 3);
\end{lstlisting}

Nazwy \texttt{xs} nie używamy w ciele funkcji do odczytywania
wartości; operacja liczenia jest czysto składniowa.

\section{Pierwsze praktyczne funkcje}

\subsection{Wypisywanie dowolnej liczby argumentów}

Najprostszy sposób wykorzystania pakietu to jego rozwinięcie
w wywołaniu innej funkcji:

\begin{lstlisting}
#include <iostream>
#include <string>

void output(int x, double y, const std::string& z) {
    std::cout << x << ' ' << y << ' ' << z << '\n';
}

template<class... Args>
void forward_to_output(Args... args) {
    output(args...);
}

int main() {
    forward_to_output(3, 2.5, std::string{"hello"});
}
\end{lstlisting}

Funkcja jest poprawna tylko dla takich wywołań, z których można
utworzyć poprawne wywołanie \texttt{output}. Sama wariadyczność
nie oznacza, że każda liczba i każdy rodzaj argumentów będą
akceptowane przez funkcję docelową.

\subsection{Przekazywanie przez wartość i przez referencję}

\begin{lstlisting}
template<class... Ts>
void by_value(Ts... xs);            // may copy/move

template<class... Ts>
void by_const_ref(const Ts&... xs); // read-only; no copies

template<class... Ts>
void by_ref(Ts&... xs);             // needs lvalues
\end{lstlisting}

Wybór postaci parametrów ma te same skutki co w zwykłych funkcjach.
W szczególności \texttt{Ts... xs} tworzy osobne parametry przez wartość,
a \texttt{const Ts\&... xs} zwykle pozwala uniknąć niepotrzebnego
kopiowania. Przypadek \texttt{Ts\&\&... xs} omówimy dokładniej
w sekcji o \emph{perfect forwarding}.

\section{Rekurencyjne przetwarzanie pakietu: metoda C++11}
\label{sec:variadic-rekurencja}

Przed C++17 często rozbijano pakiet na \emph{pierwszy argument}
i \emph{pozostałe argumenty}, a następnie wywoływano funkcję
rekurencyjnie. Potrzebny był osobny przypadek kończący rekurencję.

\begin{lstlisting}
#include <iostream>

void print_recursive() {
    std::cout << '\n';        // base case
}

template<class First, class... Rest>
void print_recursive(const First& first,
                     const Rest&... rest) {
    std::cout << first << ' ';
    print_recursive(rest...); // one argument fewer
}

int main() {
    print_recursive(1, "two", 3.0);
}
\end{lstlisting}

Przy trzech argumentach instancjonowane są kolejno wersje
z pakietami o długościach $3$, $2$, $1$ oraz przypadek końcowy.
Jest to \emph{rekurencja szablonowa}; nie należy automatycznie
utożsamiać jej z dynamiczną rekurencją wykonywaną w programie.
Kompilator może wygenerować zwykłą sekwencję operacji.

\paragraph{Kiedy stosować?}
Rekurencyjne rozkładanie pakietu nadal bywa przydatne, jeśli kolejne
kroki mają odrębne warunki lub należy budować typy rekurencyjne.
Jeżeli wykonujemy tę samą operację dla wszystkich argumentów,
w C++17 i nowszych zazwyczaj lepszym wyborem jest \emph{fold expression}.

\section{Fold expressions (C++17)}

\begin{center}\small
\begin{tabularx}{\textwidth}{@{}p{0.35\textwidth}X@{}}
\toprule
Forma złożenia & Sposób rozwinięcia \\
\midrule
\code{(... + args)} & Złożenie lewostronne: kolejne dodawania łączone od lewej. \\
\code{(args + ...)} & Złożenie prawostronne: kolejne dodawania łączone od prawej. \\
\code{(init + ... + args)} & Złożenie z wartością początkową; działa także dla pustego pakietu. \\
\code{(f(args), ...)} & Wywołuje \code{f} dla każdego argumentu od lewej do prawej dzięki operatorowi przecinka. \\
\texttt{(args \&\& ...)} & Łączy warunki z krótkim spięciem; dla pustego pakietu daje \code{true}. \\
\bottomrule
\end{tabularx}
\end{center}
\label{sec:variadic-fold}

\emph{Fold expression} redukuje pakiet wyrażeń z użyciem jednego
operatora binarnego, np. \texttt{+}, \texttt{\&\&}, \texttt{||}
lub operatora przecinka. Nie wymaga jawnego przypadku
rekurencyjnego.

\subsection{Cztery postacie składni}

Niech \texttt{xs} oznacza pakiet, $\circ$ ustalony operator,
a $z$ wartość początkową. Istnieją cztery postacie składni:

\begin{center}
\small
\begin{tabular}{|p{0.39\textwidth}|p{0.47\textwidth}|}
\hline
\textbf{Postać C++} & \textbf{Sposób łączenia dla trzech argumentów} \\
\hline
\texttt{(xs op ...)} & $x_1\circ(x_2\circ x_3)$ (prawostronne, bez inicjalizatora) \\
\hline
\texttt{(... op xs)} & $(x_1\circ x_2)\circ x_3$ (lewostronne, bez inicjalizatora) \\
\hline
\texttt{(xs op ... op z)} & $x_1\circ(x_2\circ(x_3\circ z))$ (prawostronne, z inicjalizatorem) \\
\hline
\texttt{(z op ... op xs)} & $((z\circ x_1)\circ x_2)\circ x_3$ (lewostronne, z inicjalizatorem) \\
\hline
\end{tabular}
\end{center}

W tabeli \texttt{op} jest jedynie oznaczeniem miejsca na rzeczywisty
operator, np. \texttt{+}; nie jest słowem kluczowym języka.
\emph{Nawiasy całego fold expression są obowiązkowe}.
Dwa operatory występujące w fold z inicjalizatorem muszą być takie same.

\subsection{Suma, iloczyn i warunki logiczne}

\begin{lstlisting}
#include <concepts>

template<class... Ts>
    requires (std::integral<Ts> && ...)
constexpr auto sum(Ts... xs) {
    return (0 + ... + xs);       // left fold with identity 0
}

template<class... Ts>
constexpr bool all_true(Ts... xs) {
    return (true && ... && static_cast<bool>(xs));
}

template<class... Ts>
constexpr bool any_true(Ts... xs) {
    return (false || ... || static_cast<bool>(xs));
}

static_assert(sum() == 0);
static_assert(sum(2, 3, 4) == 9);
static_assert(all_true());
static_assert(!any_true());
static_assert(all_true(1, 2, 3));
\end{lstlisting}

Wyrażenie \texttt{(0 + ... + xs)} może być używane także dla
pustego pakietu. Początkowe \texttt{0} wpływa jednak na typ wyniku
i ewentualne promocje arytmetyczne; nie jest to uniwersalna
implementacja sumy dowolnych typów liczbowych.

\subsection{Lewostronne i prawostronne składanie nie zawsze są równe}

Dla operatora odejmowania kolejność nawiasowania zmienia wynik:

\begin{lstlisting}
template<class... Ts>
constexpr auto left_sub(Ts... xs) {
    return (100 - ... - xs);
}

template<class... Ts>
constexpr auto right_sub(Ts... xs) {
    return (xs - ... - 100);
}

static_assert(left_sub(10, 3) == 87);  // (100 - 10) - 3
static_assert(right_sub(10, 3) == 107); // 10 - (3 - 100)
\end{lstlisting}

Fold określa \emph{strukturę nawiasów}, ale nie należy mylić tej
własności z ogólną gwarancją kolejności wykonywania wszystkich
podwyrażeń dla dowolnego operatora. Reguły ewaluacji wynikają
również z semantyki danego operatora.

\subsection{Pusty pakiet i element neutralny}

Dla \emph{jednoargumentowego} fold z pustym pakietem C++ dopuszcza
zaledwie trzy operatory: \texttt{\&\&} zwraca \texttt{true},
\texttt{||} zwraca \texttt{false}, a \texttt{,} daje
wyrażenie typu \texttt{void}. Pusty pakiet w
\texttt{(... + xs)} nie jest poprawny. Stosując wersję
z wartością początkową, możemy jawnie obsłużyć pusty pakiet:

\begin{lstlisting}
template<class... Ts>
constexpr auto product(Ts... xs) {
    return (1 * ... * xs);
}

static_assert(product() == 1);
static_assert(product(2, 3, 4) == 24);
\end{lstlisting}

\section{Operator przecinka i kolejność wykonywania}
\label{sec:variadic-przecinek}

Operator przecinka jest wygodny, gdy dla każdego elementu pakietu
chcemy wykonać instrukcję o efektach ubocznych. W odróżnieniu
od kolejności ewaluacji zwykłych argumentów funkcji, jego
wbudowana postać gwarantuje wykonanie lewej strony przed prawą.

\begin{lstlisting}
#include <iostream>

template<class... Ts>
void print_all(const Ts&... xs) {
    ((std::cout << xs << ' '), ...);
    std::cout << '\n';
}

int main() {
    print_all(10, 2.5, "hello");
    print_all();                // only a newline
}
\end{lstlisting}

Zewnętrzne nawiasy tworzą fold, a wewnętrzne oddzielają
wyrażenie wypisujące od operatora \texttt{,}.

\paragraph{Kolejność argumentów zwykłej funkcji.}

\begin{lstlisting}
template<class... Ts>
void sink(Ts...);

// sink(change_a(), change_b(), change_c());
\end{lstlisting}

Samo rozwinięcie \texttt{...} w liście argumentów funkcji
\emph{nie} gwarantuje, że \texttt{change\_a()} zostanie obliczone
przed \texttt{change\_b()}. W C++17 i nowszych wywołania
te nie przeplatają swoich efektów, ale ich wzajemna kolejność
nie jest określona. Gdy kolejność jest istotna, można użyć
operatora przecinka lub innej konstrukcji dającej odpowiednie
gwarancje sekwencjonowania.

\section{Pakiety typów i metaprogramowanie}
\label{sec:variadic-type-list}

Wariadyczność może dotyczyć nie tylko funkcji, lecz także
\emph{szablonów klas}. W takim przypadku pakiet przechowuje
listę typów, a instancjonowanie klasy odbywa się podczas kompilacji.

\subsection{Własna lista typów}

\begin{lstlisting}
#include <cstddef>
#include <type_traits>

template<class... Ts>
struct type_list {};

template<class List>
struct list_size;

template<class... Ts>
struct list_size<type_list<Ts...>>
    : std::integral_constant<std::size_t, sizeof...(Ts)> {};

using Types = type_list<int, double, char>;
static_assert(list_size<Types>::value == 3);
static_assert(list_size<type_list<>>::value == 0);
\end{lstlisting}

Jest to przykład \emph{częściowej specjalizacji} szablonu klasy.
Wzorzec \texttt{type\_list<Ts...>} dopasowuje dowolną listę
typów, a \texttt{sizeof...(Ts)} oblicza jej długość.

\subsection{Wydobywanie pierwszego typu}

\begin{lstlisting}
template<class List>
struct first_type;   // intentionally undefined for empty list

template<class T, class... Rest>
struct first_type<type_list<T, Rest...>> {
    using type = T;
};

using X = first_type<type_list<long, float>>::type;
static_assert(std::is_same_v<X, long>);
\end{lstlisting}

Rozkład na \texttt{T} oraz \texttt{Rest...} jest odpowiednikiem
rekurencyjnego rozdzielenia argumentów funkcji na pierwszy
i pozostałe. Dla pustej listy \texttt{type\_list<>} nazwa
\texttt{first\_type<...>::type} celowo nie istnieje.

\subsection{Czy wszystkie typy spełniają własność?}

\begin{lstlisting}
#include <type_traits>

template<class... Ts>
inline constexpr bool all_integral_v =
    (std::is_integral_v<Ts> && ...);

template<class... Ts>
inline constexpr std::size_t count_integral_v =
    (std::size_t{0} + ... +
        (std::is_integral_v<Ts> ? 1u : 0u));

static_assert(all_integral_v<int, long, char>);
static_assert(!all_integral_v<int, double>);
static_assert(count_integral_v<int, double, char> == 2);
\end{lstlisting}

Jest to częsty wzorzec: pakiet typów łączymy z cechami
typów (\emph{type traits}) i operatorami logicznymi fold.

\section{Pakiety wartości i parametry nietypowe}

Pakiet może także zawierać wartości znane w czasie kompilacji,
a nie typy. W szablonie klasy służą do tego
\emph{nietypowe parametry szablonu} (\emph{non-type template parameters}).

\begin{lstlisting}
#include <array>
#include <cstddef>

template<int... Ns>
struct int_pack {
    static constexpr std::size_t size = sizeof...(Ns);
    static constexpr int sum = (0 + ... + Ns);
};

static_assert(int_pack<>::sum == 0);
static_assert(int_pack<2, 4, 6>::sum == 12);
\end{lstlisting}

Od C++17 można zadeklarować pakiet \texttt{auto...}:

\begin{lstlisting}
template<auto... Values>
struct value_pack {
    static constexpr std::size_t size = sizeof...(Values);
};

static_assert(value_pack<42, 'x', true>::size == 3);
\end{lstlisting}

Dla \texttt{auto...} dopuszczalność konkretnej wartości jako
argumentu szablonu nadal podlega regułom parametrów nietypowych;
nie można przekazać do niego dowolnej wartości z czasu wykonania.

\section{Przekazywanie doskonałe (\textit{perfect forwarding})}
\label{sec:variadic-forwarding}

Jednym z najważniejszych zastosowań pakietów jest budowanie
\emph{funkcji pośredniczących}, które przekazują argumenty do
innej funkcji lub konstruktora, możliwie wiernie zachowując
ich kategorię wartości (\emph{value category}).

\subsection{Różnica między lvalue a rvalue}

\begin{lstlisting}
#include <iostream>

void use(int&)  { std::cout << "lvalue\n"; }
void use(int&&) { std::cout << "rvalue\n"; }

int main() {
    int x = 10;
    use(x);       // lvalue
    use(20);      // rvalue
}
\end{lstlisting}

Parametr funkcji, choćby zadeklarowany jako \texttt{T\&\&}, jest
\emph{nazwanym obiektem}; użycie jego nazwy w wyrażeniu ma
kategorię \emph{lvalue}. Dlatego samo \texttt{use(arg)}
nie zachowa kategorii pierwotnego argumentu.

\subsection{Referencje przekazujące i \texttt{std::forward}}

\begin{lstlisting}
#include <utility>

template<class T>
void wrapper(T&& arg) {
    use(std::forward<T>(arg));
}

void forwarding_example() {
    int y = 0;
    wrapper(y);  // T = int&, forwarded as lvalue
    wrapper(30); // T = int, forwarded as rvalue
}
\end{lstlisting}

\texttt{T\&\&} jest \emph{referencją przekazującą}
(ang. \emph{forwarding reference}), gdy \texttt{T} jest
odpowiednio dedukowanym, niekwalifikowanym parametrem typu.
Działa tu \emph{składanie referencji}:

\[
T\&\;\& \longrightarrow T\&,
\qquad
T\&\;\&\& \longrightarrow T\&,
\qquad
T\&\&\;\&\& \longrightarrow T\&\&.
\]

W zapisie \texttt{std::forward<T>(arg)} jawny argument
\texttt{T} jest konieczny, by odtworzyć pierwotną kategorię.
\texttt{std::forward} nie wykonuje bezwarunkowego przeniesienia:
przepuszcza lvalue jako lvalue, a rvalue jako rvalue.

\subsection{Pakiet referencji przekazujących}

\begin{lstlisting}
#include <functional>
#include <utility>

template<class F, class... Args>
decltype(auto) call(F&& f, Args&&... args) {
    return std::invoke(
        std::forward<F>(f),
        std::forward<Args>(args)...
    );
}
\end{lstlisting}

Tutaj rozwijany jest \emph{cały wzorzec}
\texttt{std::forward<Args>(args)}, a nie tylko nazwa
\texttt{args}. Każdy element używa właściwego typu z pakietu
\texttt{Args}.
\texttt{std::invoke} potrafi wywoływać zwykłe funkcje,
obiekty funkcyjne i odpowiednio obsłużone wskaźniki do składowych.
\texttt{decltype(auto)} zachowuje typ zwracany przez docelowe
wywołanie, w tym referencję, jeżeli docelowa funkcja ją zwraca.

\begin{lstlisting}
#include <string>

std::string join(std::string a, std::string b) {
    return a + b;
}

auto result = call(join, std::string{"ab"}, "cd");
// result == "abcd"
\end{lstlisting}

\subsection{Czego nie robić: \texttt{std::move} i nazwy parametrów}

\begin{lstlisting}
template<class T>
void wrong_1(T&& arg) {
    use(arg);              // arg is always an lvalue here
}

template<class T>
void wrong_2(T&& arg) {
    use(std::move(arg));   // forces an rvalue even for lvalues
}
\end{lstlisting}

\texttt{std::move} to rzutowanie umożliwiające potraktowanie
obiektu jak rvalue, a nie polecenie ``zachowaj kategorię''.
Stosuj je, gdy rzeczywiście przekazujesz obiekt do operacji,
która może pobrać z niego zasób. W uniwersalnym wrapperze
używaj \texttt{std::forward}.

\paragraph{Ograniczenia żywotności.}
Referencje przekazujące \emph{nie przedłużają automatycznie}
żywotności obiektu do czasu dowolnego późniejszego wykonania.
Przechowanie referencji do argumentu tymczasowego w strukturze,
która przeżyje wywołanie, może prowadzić do wskaźników lub
referencji wiszących. Ten sam problem dotyczy między innymi
\texttt{std::forward\_as\_tuple}.

\section{Tworzenie obiektów z dowolnych argumentów}
\label{sec:variadic-emplace}

\subsection{Własna fabryka obiektów}

\begin{lstlisting}
#include <memory>
#include <utility>

template<class T, class... Args>
std::unique_ptr<T> create(Args&&... args) {
    return std::make_unique<T>(
        std::forward<Args>(args)...
    );
}

struct Point {
    int x, y;
    Point(int a, int b) : x(a), y(b) {}
};

auto point = create<Point>(4, 7);
\end{lstlisting}

\texttt{std::make\_unique} już dostarcza taką funkcjonalność,
więc \texttt{create} ma tu charakter dydaktyczny. Pokazuje
jednak wzorzec używany w fabrykach oraz funkcjach
\texttt{emplace} i \texttt{make}.

\subsection{Tworzenie kontenera bez wymogu kopiowalności}

\begin{lstlisting}
#include <concepts>
#include <string>
#include <utility>
#include <vector>

template<class T, class... Args>
    requires (std::constructible_from<T, Args&&> && ...)
std::vector<T> make_vector(Args&&... args) {
    std::vector<T> result;
    result.reserve(sizeof...(Args));
    (result.emplace_back(std::forward<Args>(args)), ...);
    return result;
}

auto words = make_vector<std::string>("A", "BB", "CCC");
\end{lstlisting}

Operator przecinka daje tu jednoznaczną kolejność wywołań
\texttt{emplace\_back}. Funkcja przyjmuje także pusty pakiet
argumentów. W przeciwieństwie do budowania wektora z listy
inicjalizacyjnej pozwala łatwo przekazywać do kolejnych
\texttt{emplace\_back} argumenty przenoszalne, ale niekopiowalne.
Dla elementów konstruowanych z więcej niż jednego parametru
wymagałaby innego interfejsu grupowania argumentów.

\section{\texttt{std::tuple}, \texttt{std::apply} i \texttt{std::index\_sequence}}
\label{sec:variadic-tuple}

\subsection{Heterogeniczna krotka \texttt{std::tuple}}

\texttt{std::tuple<Ts...>} to kontener o strukturze opisanej
listą typów znaną w czasie kompilacji. Każda pozycja może
mieć inny typ:

\begin{lstlisting}
#include <tuple>
#include <string>

std::tuple<int, std::string, double> t{7, "Ada", 3.5};

int n = std::get<0>(t);
const std::string& name = std::get<1>(t);
\end{lstlisting}

W przeciwieństwie do \texttt{std::vector<T>}, krotka nie
przechowuje dowolnej liczby elementów jednego typu, lecz
\emph{ustaloną na etapie kompilacji} listę typów.

\subsection{Rozwinięcie krotki do argumentów funkcji}

Funkcja \texttt{std::apply} (C++17) wywołuje funkcję,
przekazując elementy obiektu \emph{tuple-like} jako osobne
argumenty:

\begin{lstlisting}
#include <iostream>
#include <tuple>

void apply_example() {
    std::apply([](const auto&... xs) {
        ((std::cout << xs << ' '), ...);
        std::cout << '\n';
    }, t);
}
\end{lstlisting}

\texttt{std::apply} oszczędza ręcznego konstruowania
sekwencji indeksów, ale warto wiedzieć, jak taki mechanizm
można zaimplementować.

\subsection{Pakiety indeksów: \texttt{std::index\_sequence}}

\begin{lstlisting}
#include <cstddef>
#include <iostream>
#include <tuple>
#include <utility>

template<class Tuple, std::size_t... I>
void print_tuple_impl(const Tuple& t,
                      std::index_sequence<I...>) {
    ((std::cout << std::get<I>(t) << ' '), ...);
    std::cout << '\n';
}

template<class... Ts>
void print_tuple(const std::tuple<Ts...>& t) {
    print_tuple_impl(t, std::index_sequence_for<Ts...>{});
}

void tuple_example() {
    print_tuple(std::make_tuple(1, "abc", 3.0));
}
\end{lstlisting}

\texttt{std::index\_sequence<0,1,2>} jest specjalizacją
szablonu przyjmującego pakiet indeksów całkowitych.
\texttt{std::index\_sequence\_for<Ts...>} generuje indeksy
$0,1,\ldots,n-1$ dla pakietu długości $n$.
Wzorzec \texttt{std::get<I>(t)...} rozwija wyrażenie
osobno dla każdego indeksu.

\subsection{Równoległe rozwijanie pakietów}

W jednym wzorcu można użyć kilku pakietów. Muszą mieć
\emph{tę samą długość}, jeśli są rozwijane równolegle:

\begin{lstlisting}
#include <tuple>
#include <type_traits>
#include <utility>

template<class A, class B>
struct zipped_types;

template<class... Ts, class... Us>
struct zipped_types<std::tuple<Ts...>, std::tuple<Us...>> {
    static_assert(sizeof...(Ts) == sizeof...(Us));
    using type = std::tuple<std::pair<Ts, Us>...>;
};

using Z = zipped_types<
    std::tuple<int, char>,
    std::tuple<double, bool>
>::type;

static_assert(std::is_same_v<
    Z, std::tuple<std::pair<int, double>,
                  std::pair<char, bool>>>);
\end{lstlisting}

Elementy \texttt{Ts} i \texttt{Us} są łączone pozycjami:
$\bigl(T_1,U_1\bigr),\ldots,\bigl(T_n,U_n\bigr)$.
Nie jest to iloczyn kartezjański pakietów.

\section{Dziedziczenie wariadyczne i \texttt{std::visit}}
\label{sec:variadic-overloaded}

Pakiet \texttt{Bases...} można rozwinąć na liście baz.
Z kolei zapis \texttt{using Bases::operator()...}
wprowadza ich operatory wywołania do jednej klasy.
W ten sposób otrzymujemy zestaw przeciążeń.

\begin{lstlisting}
#include <iostream>
#include <string>
#include <variant>

template<class... Fs>
struct overloaded : Fs... {
    using Fs::operator()...;
};

template<class... Fs>
overloaded(Fs...) -> overloaded<Fs...>; // C++17 deduction guide

int main() {
    std::variant<int, std::string> value = "hello";

    std::visit(overloaded{
        [](int n) { std::cout << "int: " << n << '\n'; },
        [](const std::string& s) {
            std::cout << "string: " << s << '\n';
        }
    }, value);
}
\end{lstlisting}

Jest to standardowy idiom C++17: każda lambda definiuje
osobny \texttt{operator()}, a dziedziczenie wariadyczne
łączy je w jeden obiekt wywoływalny. Gdy stosujemy ten wzorzec
w wielu miejscach, warto umieścić \texttt{overloaded} w
niewielkim nagłówku narzędziowym.

\section{Ograniczanie pakietów za pomocą \texttt{concepts} (C++20)}
\label{sec:variadic-concepts}

Nie każda operacja ma sens dla każdej możliwej kombinacji
typów. W C++20 można czytelnie określać wymagania dla
\emph{wszystkich} elementów pakietu.

\begin{lstlisting}
#include <concepts>
#include <type_traits>

template<class... Ts>
    requires (std::integral<Ts> && ...)
constexpr bool all_nonnegative(Ts... xs) {
    return ((xs >= 0) && ...);
}

static_assert(all_nonnegative(0, 2, 7));
static_assert(!all_nonnegative(0, -1, 7));

// all_nonnegative(3.14); // rejected by constraints
\end{lstlisting}

Możemy również ograniczać pakiet w odniesieniu do
konkretnego typu docelowego:

\begin{lstlisting}
#include <concepts>

template<class T, class... Args>
concept ConstructibleWith =
    std::constructible_from<T, Args...>;

static_assert(ConstructibleWith<std::string, const char*>);
static_assert(!ConstructibleWith<int, std::string>);
\end{lstlisting}

Pierwszy przykład stosuje fold nad \emph{warunkami dla pojedynczych
typów}; drugi sprawdza \emph{jedną własność całej sekwencji
typów}, tj. czy odpowiedni konstruktor przyjmuje argumenty
\texttt{Args...} w tej kolejności. Tych dwóch mechanizmów
nie należy utożsamiać.

\subsection{\texttt{if constexpr} i struktura pakietu}

\begin{lstlisting}
#include <iostream>

template<class... Ts>
void report(const Ts&... xs) {
    if constexpr (sizeof...(Ts) == 0) {
        std::cout << "empty\n";
    } else {
        ((std::cout << xs << ' '), ...);
        std::cout << '\n';
    }
}
\end{lstlisting}

\texttt{if constexpr} pozwala w zależności od długości pakietu
wybierać gałąź w czasie kompilacji. W dobrze sformułowanym
szablonie odrzucona gałąź nie jest instancjonowana dla danej
specjalizacji. To użyteczne przy przypadkach brzegowych,
które nie mają sensu dla pustego pakietu.

\section{Ograniczenia, pułapki i diagnostyka}
\label{sec:variadic-pulapki}

\subsection{Pakiet typów nie jest pakietem argumentów funkcji}

\begin{lstlisting}
template<class... Ts>
void f(Ts... args) {
    // sizeof...(Ts): number of types
    // sizeof...(args): number of function parameters
    // args[0];  // ERROR: args is not an array
}
\end{lstlisting}

Do wybierania argumentu na ustalonej pozycji można użyć
\texttt{std::tuple} i \texttt{std::get} albo wywoływać funkcję
przez \texttt{std::apply}. Sam pakiet \texttt{args...}
nie jest indeksowanym kontenerem.

\subsection{Brak zgodności typów po rozwinięciu}

\begin{lstlisting}
template<class... Ts>
auto add(Ts... xs) {
    return (0 + ... + xs);
}

// add(1, std::string{"x"}); // ERROR: unsuitable operator+
\end{lstlisting}

Pakiet nie omija systemu typów. Każde wyrażenie powstałe po
jego rozwinięciu musi być poprawne, podobnie jak zwykły kod
C++. Jeżeli szablon ma określone wymagania, najlepiej
zapisać je za pomocą \texttt{requires}.

\subsection{Niewłaściwe przekazywanie kategorii wartości}

\begin{lstlisting}
template<class F, class... Args>
void bad_call(F&& f, Args&&... args) {
    std::forward<F>(f)(args...); // loses rvalue categories
}

// Correct:
// std::forward<F>(f)(std::forward<Args>(args)...);
\end{lstlisting}

Dla typów niekopiowalnych taki błąd często ujawnia się
natychmiast podczas kompilacji. Dla typów kopiowalnych
może spowodować jedynie dodatkowe kopie albo zmianę
wybranego przeciążenia.

\subsection{Wnioskowanie typów i \texttt{\{...\}}}

\begin{lstlisting}
#include <initializer_list>

template<class... Args>
void foo(Args&&...) {}

void braced_example() {
    // foo({1, 2, 3}); // ERROR: cannot deduce Args
    foo(std::initializer_list<int>{1, 2, 3}); // OK
}
\end{lstlisting}

Sama lista ujęta w nawiasy klamrowe \texttt{\{...\}} nie ma
typu, który dowolny pakiet \texttt{Args\&\&...} potrafiłby
wydedukować. Potrzebny jest określony typ parametru,
np. \texttt{std::initializer\_list}\allowbreak\texttt{<int>}, lub jawnie
utworzony obiekt.

\subsection{Ograniczenia położenia pakietu}

W \emph{głównej definicji szablonu klasy} pakiet parametrów
szablonu musi być ostatni. Dla szablonów funkcji istnieją
bardziej elastyczne reguły, ale kolejne parametry muszą
być dedukowalne lub mieć argumenty domyślne, stosownie do
kontekstu:

\begin{lstlisting}
// ERROR for a primary class template:
// template<class... Ts, class U> struct invalid;

template<class U, class... Ts>
struct valid {};
\end{lstlisting}

\subsection{Skomplikowany kod zamiast prostego rozwiązania}

Szablon wariadyczny to narzędzie, a nie obowiązkowy styl.
Jeżeli argumenty mają \emph{jednolity typ i ich liczba
zmienia się w czasie wykonania}, lepszy bywa
\texttt{std::span<T>}, \texttt{std::vector<T>} lub para
iteratorów. Pakiet ma liczbę elementów i strukturę typów
ustalaną na etapie kompilacji; nie zastępuje kontenera.

\section{Variadic templates a funkcje wariadyczne w stylu C}

Historyczne funkcje z deklaracją \texttt{void f(int, ...)}
korzystają z mechanizmu argumentów wariadycznych języka C,
a nie z pakietu parametrów szablonu C++:

\begin{lstlisting}
#include <cstdio>

int main() {
    std::printf("%d %f\n", 7, 2.5); // C-style variadic call
}
\end{lstlisting}

W formacie C zgodność argumentów ze specyfikatorami
formatowania jest obowiązkiem programisty (niektóre
kompilatory potrafią wykrywać część pomyłek). W szablonach
wariadycznych typy argumentów pozostają dostępne
kompilatorowi, dzięki czemu można wymuszać ograniczenia
i stosować przeciążanie oraz \texttt{std::forward}.
Nie należy jednak wyciągać z tego wniosku, że dowolny
szablon automatycznie bezbłędnie formatuje argumenty.

\section{Pełny przykład: prosty rejestrator zdarzeń}
\label{sec:variadic-projekt}

Zastosujmy omówione mechanizmy do funkcji tworzącej tekstowy
komunikat z dowolnej liczby argumentów możliwych do
wypisania do strumienia. Zamiast ręcznie doklejać spacje
w wielu przeciążeniach, użyjemy fold nad operatorem przecinka.

\begin{lstlisting}
#include <iostream>
#include <sstream>
#include <string>

// Each argument must support insertion into std::ostream.
template<class... Args>
std::string make_message(const Args&... args) {
    std::ostringstream out;
    ((out << args), ...);
    return out.str();
}

template<class... Args>
void log_line(const Args&... args) {
    std::cout << "[LOG] " << make_message(args...) << '\n';
}

int main() {
    log_line("count=", 3, ", ratio=", 0.75);
    log_line("finished");
    log_line();
}
\end{lstlisting}

\paragraph{Analiza.}
Dla argumentów \texttt{("count=",3,", ratio=",0.75)}
pakiet zawiera cztery elementy. Wyrażenie
\texttt{((out << args), ...)} wstawia je kolejno do strumienia,
a \texttt{make\_message} zwraca powstały napis. Dla
pustego pakietu fold z operatorem przecinka jest poprawny,
a wynikowy napis jest pusty. W wersji produkcyjnej należałoby
rozważyć wymagania dotyczące formatowania, zachowanie w
obliczu wyjątków, współbieżność oraz ewentualne skorzystanie
z narzędzia formatującego wspieranego przez używany kompilator.

\section{Zestawienie: kiedy używać poszczególnych mechanizmów?}

\begin{center}
\small
\begin{tabular}{|p{0.40\textwidth}|p{0.46\textwidth}|}
\hline
\textbf{Cel} & \textbf{Typowe rozwiązanie} \\
\hline
Dowolna liczba argumentów o różnych typach & \texttt{template<class... Ts>} i \texttt{Ts...} \\
\hline
Wykonanie tej samej operacji dla każdego argumentu & Fold z operatorem \texttt{,} (C++17) \\
\hline
Suma, iloczyn lub wspólny warunek & Fold z \texttt{+}, \texttt{*}, \texttt{\&\&} lub \texttt{||} \\
\hline
Przekazanie argumentów dalej z zachowaniem kategorii & \texttt{Args\&\&...}, \texttt{std::forward<Args>(args)...} \\
\hline
Zapis heterogenicznych wartości & \texttt{std::tuple<Ts...>} \\
\hline
Wywołanie funkcji elementami krotki & \texttt{std::apply} (C++17) \\
\hline
Rozwinięcie krotki przy pomocy indeksów & \texttt{std::index\_sequence} \\
\hline
Lista typów do obliczeń podczas kompilacji & Wariadyczny szablon klasy i specjalizacje \\
\hline
Ograniczenie do odpowiednich typów & Fold z koncepcjami i \texttt{requires} (C++20) \\
\hline
Zmienna długość listy jednolitych danych w czasie wykonania & \texttt{std::vector} lub \texttt{std::span}, nie pakiet \\
\hline
\end{tabular}
\end{center}

\section{Podsumowanie i zadania do samodzielnego sprawdzenia}

\begin{enumerate}
    \item \texttt{class... Ts} deklaruje \emph{pakiet typów},
    a \texttt{Ts... xs} wprowadza \emph{pakiet parametrów funkcji}.
    \item \texttt{sizeof...(Ts)} oblicza liczbę elementów pakietu
    podczas kompilacji.
    \item \texttt{pattern...} rozwija cały \emph{wzorzec},
    nie tylko pojedynczy identyfikator.
    \item Fold expressions (od C++17) zastępują wiele prostych
    rekurencyjnych implementacji z C++11.
    \item Dla pustego pakietu istotne są element neutralny i
    reguły operatora; nie każda forma fold jest poprawna.
    \item W wrapperach zachowuj kategorię wartości argumentów,
    używając referencji przekazujących i \texttt{std::forward}.
    \item \texttt{std::tuple}, \texttt{std::apply} i
    \texttt{std::index\_sequence} współpracują z pakietami,
    umożliwiając operacje na heterogenicznych danych.
    \item W C++20 używaj \texttt{concepts}/\texttt{requires},
    aby błędne wywołania odrzucać na granicy interfejsu.
\end{enumerate}

\paragraph{Pytania kontrolne.}
Spróbuj przewidzieć wyniki poniższych wyrażeń bez uruchamiania
programu. Następnie zweryfikuj je za pomocą \texttt{static\_assert}:

\begin{lstlisting}
// static_assert(sizeof...(int) == 1); // ERROR: int is not a pack

// For the functions defined above:
static_assert(product() == 1);
static_assert(left_sub(10, 3) == 87);
static_assert(right_sub(10, 3) == 107);
static_assert(count_integral_v<int, double, char> == 2);
\end{lstlisting}

Pierwsza linia jest \emph{celowo niepoprawna}: operator
\texttt{sizeof...} wymaga nazwy pakietu, a nie zwykłego typu.

% Wykorzystane punkty odniesienia (bez kopiowania tresci):
% https://en.cppreference.com/w/cpp/language/parameter_pack
% https://en.cppreference.com/w/cpp/language/fold
% https://en.cppreference.com/w/cpp/utility/forward
% https://en.cppreference.com/w/cpp/utility/apply
% https://en.cppreference.com/w/cpp/utility/integer_sequence


\chapter[Variadic templates w C++]{Szablony o zmiennej liczbie argumentów (\textit{variadic templates})}
\label{chap:variadic-templates-cpp}

\section{Idea i zastosowania}

Zwykły szablon funkcji przyjmuje z góry ustaloną liczbę parametrów
szablonu. \emph{Variadic template} (szablon o zmiennej liczbie parametrów)
może natomiast przyjmować \emph{dowolną liczbę} argumentów szablonu,
również zero. Mechanizmem, który to umożliwia, jest
\emph{pakiet parametrów} (ang. \emph{parameter pack}).

\begin{lstlisting}
template<class T>
void single(const T&) {}        // one type

template<class... Ts>
void many(const Ts&...) {}      // zero or more types

void demo() {
    many();                     // zero arguments
    many(10);                   // one argument
    many(10, 3.5, "abc");      // three arguments
}
\end{lstlisting}

Pakiety parametrów są ważną częścią nowoczesnego C++.
Wykorzystują je m.in. klasy \texttt{std::tuple} i \texttt{std::variant}
oraz fabryki \texttt{std::make\_unique} i \texttt{std::make\_shared}.
Nie trzeba dzięki nim pisać przeciążeń dla każdej
liczby argumentów.

\paragraph{Wymagane wersje języka.}
Pakiety parametrów i perfect forwarding pojawiły się w C++11.
\emph{Fold expressions} dodano w C++17, a koncepcje
(\texttt{concepts}) i klauzulę \texttt{requires} w C++20.
Dalej zakładamy C++20.

\section{Trzy kropki: deklaracja pakietu a rozwinięcie}

\begin{center}\small
\begin{tabularx}{\textwidth}{@{}p{0.35\textwidth}X@{}}
\toprule
Zapis & Działanie \\
\midrule
\code{typename... Ts} & Deklaruje pakiet parametrów typowych szablonu. \\
\code{Ts...} & Rozwija pakiet typów w miejscu, które dopuszcza taki wzorzec. \\
\code{args...} & Rozwija pakiet argumentów, tworząc listę wyrażeń. \\
\code{sizeof...(Ts)} & Podaje liczbę elementów pakietu w czasie kompilacji. \\
\code{std::forward<Ts>(args)...} & Przekazuje każdy argument z zachowaniem jego kategorii wartości. \\
\bottomrule
\end{tabularx}
\end{center}
\label{sec:variadic-trzy-kropki}

W różnych miejscach token \texttt{...} oznacza różne operacje.
Rozróżnienie to stanowi podstawę poprawnego używania szablonów
wariadycznych.

\begin{lstlisting}
template<class... Ts>             // declare type pack Ts
void f(Ts... xs) {                // declare function parameter pack xs
    // ...
}
\end{lstlisting}

Deklaracja \texttt{class... Ts} wprowadza pakiet \emph{typów}.
Zapis \texttt{Ts... xs} rozwija pakiet typów i wprowadza odpowiadający
mu pakiet parametrów funkcji. Wywołanie
\texttt{f(1, 2.0, 'x')} tworzy w uproszczeniu funkcję o parametrach
\texttt{(int, double, char)}.

\subsection{Rozwinięcie pakietu (\textit{pack expansion})}

Rozwinięcie składa się ze \emph{wzorca} (\emph{pattern}) oraz następującego
po nim \texttt{...}. Dla każdego elementu pakietu kompilator podstawia
kolejny element do całego wzorca:

\begin{lstlisting}
template<class... Ts>
void demo(Ts... xs) {
    g(xs...);                // g(x1, x2, ..., xn)
    h((xs * 2)...);          // h(x1*2, x2*2, ..., xn*2)
}
\end{lstlisting}

W drugim wywołaniu wzorcem jest \texttt{(xs * 2)}, a \emph{nie}
pojedyncza zmienna \texttt{xs}. Ważne: \texttt{xs...} nie jest
kontenerem ani tablicą; nie można po nim bezpośrednio iterować zwykłą
pętlą \texttt{for}. To składnia kompilatora służąca do generowania
listy wyrażeń.

\paragraph{Pakiet typów a pakiet wartości.}
W przykładzie \texttt{Ts} reprezentuje typy, natomiast
\texttt{xs} to konkretne parametry funkcji. Jeden typ może wystąpić
w pakiecie wielokrotnie, a każdy element pakietu wartości zachowuje
swój właściwy typ.

\subsection{Liczba elementów: \texttt{sizeof...}}

Operator \texttt{sizeof...} daje liczbę elementów pakietu jako
stałą czasu kompilacji. Nie należy go mylić z operatorem
\texttt{sizeof} mierzącym rozmiar typu lub obiektu.

\begin{lstlisting}
#include <cstddef>

template<class... Ts>
constexpr std::size_t count(Ts&&... xs) {
    static_assert(sizeof...(Ts) == sizeof...(xs));
    return sizeof...(xs);
}

static_assert(count() == 0);
static_assert(count(2, 3.5, 'a') == 3);
\end{lstlisting}

Nazwy \texttt{xs} nie używamy w ciele funkcji do odczytywania
wartości; operacja liczenia jest czysto składniowa.

\section{Pierwsze praktyczne funkcje}

\subsection{Wypisywanie dowolnej liczby argumentów}

Najprostszy sposób wykorzystania pakietu to jego rozwinięcie
w wywołaniu innej funkcji:

\begin{lstlisting}
#include <iostream>
#include <string>

void output(int x, double y, const std::string& z) {
    std::cout << x << ' ' << y << ' ' << z << '\n';
}

template<class... Args>
void forward_to_output(Args... args) {
    output(args...);
}

int main() {
    forward_to_output(3, 2.5, std::string{"hello"});
}
\end{lstlisting}

Funkcja jest poprawna tylko dla takich wywołań, z których można
utworzyć poprawne wywołanie \texttt{output}. Sama wariadyczność
nie oznacza, że każda liczba i każdy rodzaj argumentów będą
akceptowane przez funkcję docelową.

\subsection{Przekazywanie przez wartość i przez referencję}

\begin{lstlisting}
template<class... Ts>
void by_value(Ts... xs);            // may copy/move

template<class... Ts>
void by_const_ref(const Ts&... xs); // read-only; no copies

template<class... Ts>
void by_ref(Ts&... xs);             // needs lvalues
\end{lstlisting}

Wybór postaci parametrów ma te same skutki co w zwykłych funkcjach.
W szczególności \texttt{Ts... xs} tworzy osobne parametry przez wartość,
a \texttt{const Ts\&... xs} zwykle pozwala uniknąć niepotrzebnego
kopiowania. Przypadek \texttt{Ts\&\&... xs} omówimy dokładniej
w sekcji o \emph{perfect forwarding}.

\section{Rekurencyjne przetwarzanie pakietu: metoda C++11}
\label{sec:variadic-rekurencja}

Przed C++17 często rozbijano pakiet na \emph{pierwszy argument}
i \emph{pozostałe argumenty}, a następnie wywoływano funkcję
rekurencyjnie. Potrzebny był osobny przypadek kończący rekurencję.

\begin{lstlisting}
#include <iostream>

void print_recursive() {
    std::cout << '\n';        // base case
}

template<class First, class... Rest>
void print_recursive(const First& first,
                     const Rest&... rest) {
    std::cout << first << ' ';
    print_recursive(rest...); // one argument fewer
}

int main() {
    print_recursive(1, "two", 3.0);
}
\end{lstlisting}

Przy trzech argumentach instancjonowane są kolejno wersje
z pakietami o długościach $3$, $2$, $1$ oraz przypadek końcowy.
Jest to \emph{rekurencja szablonowa}; nie należy automatycznie
utożsamiać jej z dynamiczną rekurencją wykonywaną w programie.
Kompilator może wygenerować zwykłą sekwencję operacji.

\paragraph{Kiedy stosować?}
Rekurencyjne rozkładanie pakietu nadal bywa przydatne, jeśli kolejne
kroki mają odrębne warunki lub należy budować typy rekurencyjne.
Jeżeli wykonujemy tę samą operację dla wszystkich argumentów,
w C++17 i nowszych zazwyczaj lepszym wyborem jest \emph{fold expression}.

\section{Fold expressions (C++17)}

\begin{center}\small
\begin{tabularx}{\textwidth}{@{}p{0.35\textwidth}X@{}}
\toprule
Forma złożenia & Sposób rozwinięcia \\
\midrule
\code{(... + args)} & Złożenie lewostronne: kolejne dodawania łączone od lewej. \\
\code{(args + ...)} & Złożenie prawostronne: kolejne dodawania łączone od prawej. \\
\code{(init + ... + args)} & Złożenie z wartością początkową; działa także dla pustego pakietu. \\
\code{(f(args), ...)} & Wywołuje \code{f} dla każdego argumentu od lewej do prawej dzięki operatorowi przecinka. \\
\texttt{(args \&\& ...)} & Łączy warunki z krótkim spięciem; dla pustego pakietu daje \code{true}. \\
\bottomrule
\end{tabularx}
\end{center}
\label{sec:variadic-fold}

\emph{Fold expression} redukuje pakiet wyrażeń z użyciem jednego
operatora binarnego, np. \texttt{+}, \texttt{\&\&}, \texttt{||}
lub operatora przecinka. Nie wymaga jawnego przypadku
rekurencyjnego.

\subsection{Cztery postacie składni}

Niech \texttt{xs} oznacza pakiet, $\circ$ ustalony operator,
a $z$ wartość początkową. Istnieją cztery postacie składni:

\begin{center}
\small
\begin{tabular}{|p{0.39\textwidth}|p{0.47\textwidth}|}
\hline
\textbf{Postać C++} & \textbf{Sposób łączenia dla trzech argumentów} \\
\hline
\texttt{(xs op ...)} & $x_1\circ(x_2\circ x_3)$ (prawostronne, bez inicjalizatora) \\
\hline
\texttt{(... op xs)} & $(x_1\circ x_2)\circ x_3$ (lewostronne, bez inicjalizatora) \\
\hline
\texttt{(xs op ... op z)} & $x_1\circ(x_2\circ(x_3\circ z))$ (prawostronne, z inicjalizatorem) \\
\hline
\texttt{(z op ... op xs)} & $((z\circ x_1)\circ x_2)\circ x_3$ (lewostronne, z inicjalizatorem) \\
\hline
\end{tabular}
\end{center}

W tabeli \texttt{op} jest jedynie oznaczeniem miejsca na rzeczywisty
operator, np. \texttt{+}; nie jest słowem kluczowym języka.
\emph{Nawiasy całego fold expression są obowiązkowe}.
Dwa operatory występujące w fold z inicjalizatorem muszą być takie same.

\subsection{Suma, iloczyn i warunki logiczne}

\begin{lstlisting}
#include <concepts>

template<class... Ts>
    requires (std::integral<Ts> && ...)
constexpr auto sum(Ts... xs) {
    return (0 + ... + xs);       // left fold with identity 0
}

template<class... Ts>
constexpr bool all_true(Ts... xs) {
    return (true && ... && static_cast<bool>(xs));
}

template<class... Ts>
constexpr bool any_true(Ts... xs) {
    return (false || ... || static_cast<bool>(xs));
}

static_assert(sum() == 0);
static_assert(sum(2, 3, 4) == 9);
static_assert(all_true());
static_assert(!any_true());
static_assert(all_true(1, 2, 3));
\end{lstlisting}

Wyrażenie \texttt{(0 + ... + xs)} może być używane także dla
pustego pakietu. Początkowe \texttt{0} wpływa jednak na typ wyniku
i ewentualne promocje arytmetyczne; nie jest to uniwersalna
implementacja sumy dowolnych typów liczbowych.

\subsection{Lewostronne i prawostronne składanie nie zawsze są równe}

Dla operatora odejmowania kolejność nawiasowania zmienia wynik:

\begin{lstlisting}
template<class... Ts>
constexpr auto left_sub(Ts... xs) {
    return (100 - ... - xs);
}

template<class... Ts>
constexpr auto right_sub(Ts... xs) {
    return (xs - ... - 100);
}

static_assert(left_sub(10, 3) == 87);  // (100 - 10) - 3
static_assert(right_sub(10, 3) == 107); // 10 - (3 - 100)
\end{lstlisting}

Fold określa \emph{strukturę nawiasów}, ale nie należy mylić tej
własności z ogólną gwarancją kolejności wykonywania wszystkich
podwyrażeń dla dowolnego operatora. Reguły ewaluacji wynikają
również z semantyki danego operatora.

\subsection{Pusty pakiet i element neutralny}

Dla \emph{jednoargumentowego} fold z pustym pakietem C++ dopuszcza
zaledwie trzy operatory: \texttt{\&\&} zwraca \texttt{true},
\texttt{||} zwraca \texttt{false}, a \texttt{,} daje
wyrażenie typu \texttt{void}. Pusty pakiet w
\texttt{(... + xs)} nie jest poprawny. Stosując wersję
z wartością początkową, możemy jawnie obsłużyć pusty pakiet:

\begin{lstlisting}
template<class... Ts>
constexpr auto product(Ts... xs) {
    return (1 * ... * xs);
}

static_assert(product() == 1);
static_assert(product(2, 3, 4) == 24);
\end{lstlisting}

\section{Operator przecinka i kolejność wykonywania}
\label{sec:variadic-przecinek}

Operator przecinka jest wygodny, gdy dla każdego elementu pakietu
chcemy wykonać instrukcję o efektach ubocznych. W odróżnieniu
od kolejności ewaluacji zwykłych argumentów funkcji, jego
wbudowana postać gwarantuje wykonanie lewej strony przed prawą.

\begin{lstlisting}
#include <iostream>

template<class... Ts>
void print_all(const Ts&... xs) {
    ((std::cout << xs << ' '), ...);
    std::cout << '\n';
}

int main() {
    print_all(10, 2.5, "hello");
    print_all();                // only a newline
}
\end{lstlisting}

Zewnętrzne nawiasy tworzą fold, a wewnętrzne oddzielają
wyrażenie wypisujące od operatora \texttt{,}.

\paragraph{Kolejność argumentów zwykłej funkcji.}

\begin{lstlisting}
template<class... Ts>
void sink(Ts...);

// sink(change_a(), change_b(), change_c());
\end{lstlisting}

Samo rozwinięcie \texttt{...} w liście argumentów funkcji
\emph{nie} gwarantuje, że \texttt{change\_a()} zostanie obliczone
przed \texttt{change\_b()}. W C++17 i nowszych wywołania
te nie przeplatają swoich efektów, ale ich wzajemna kolejność
nie jest określona. Gdy kolejność jest istotna, można użyć
operatora przecinka lub innej konstrukcji dającej odpowiednie
gwarancje sekwencjonowania.

\section{Pakiety typów i metaprogramowanie}
\label{sec:variadic-type-list}

Wariadyczność może dotyczyć nie tylko funkcji, lecz także
\emph{szablonów klas}. W takim przypadku pakiet przechowuje
listę typów, a instancjonowanie klasy odbywa się podczas kompilacji.

\subsection{Własna lista typów}

\begin{lstlisting}
#include <cstddef>
#include <type_traits>

template<class... Ts>
struct type_list {};

template<class List>
struct list_size;

template<class... Ts>
struct list_size<type_list<Ts...>>
    : std::integral_constant<std::size_t, sizeof...(Ts)> {};

using Types = type_list<int, double, char>;
static_assert(list_size<Types>::value == 3);
static_assert(list_size<type_list<>>::value == 0);
\end{lstlisting}

Jest to przykład \emph{częściowej specjalizacji} szablonu klasy.
Wzorzec \texttt{type\_list<Ts...>} dopasowuje dowolną listę
typów, a \texttt{sizeof...(Ts)} oblicza jej długość.

\subsection{Wydobywanie pierwszego typu}

\begin{lstlisting}
template<class List>
struct first_type;   // intentionally undefined for empty list

template<class T, class... Rest>
struct first_type<type_list<T, Rest...>> {
    using type = T;
};

using X = first_type<type_list<long, float>>::type;
static_assert(std::is_same_v<X, long>);
\end{lstlisting}

Rozkład na \texttt{T} oraz \texttt{Rest...} jest odpowiednikiem
rekurencyjnego rozdzielenia argumentów funkcji na pierwszy
i pozostałe. Dla pustej listy \texttt{type\_list<>} nazwa
\texttt{first\_type<...>::type} celowo nie istnieje.

\subsection{Czy wszystkie typy spełniają własność?}

\begin{lstlisting}
#include <type_traits>

template<class... Ts>
inline constexpr bool all_integral_v =
    (std::is_integral_v<Ts> && ...);

template<class... Ts>
inline constexpr std::size_t count_integral_v =
    (std::size_t{0} + ... +
        (std::is_integral_v<Ts> ? 1u : 0u));

static_assert(all_integral_v<int, long, char>);
static_assert(!all_integral_v<int, double>);
static_assert(count_integral_v<int, double, char> == 2);
\end{lstlisting}

Jest to częsty wzorzec: pakiet typów łączymy z cechami
typów (\emph{type traits}) i operatorami logicznymi fold.

\section{Pakiety wartości i parametry nietypowe}

Pakiet może także zawierać wartości znane w czasie kompilacji,
a nie typy. W szablonie klasy służą do tego
\emph{nietypowe parametry szablonu} (\emph{non-type template parameters}).

\begin{lstlisting}
#include <array>
#include <cstddef>

template<int... Ns>
struct int_pack {
    static constexpr std::size_t size = sizeof...(Ns);
    static constexpr int sum = (0 + ... + Ns);
};

static_assert(int_pack<>::sum == 0);
static_assert(int_pack<2, 4, 6>::sum == 12);
\end{lstlisting}

Od C++17 można zadeklarować pakiet \texttt{auto...}:

\begin{lstlisting}
template<auto... Values>
struct value_pack {
    static constexpr std::size_t size = sizeof...(Values);
};

static_assert(value_pack<42, 'x', true>::size == 3);
\end{lstlisting}

Dla \texttt{auto...} dopuszczalność konkretnej wartości jako
argumentu szablonu nadal podlega regułom parametrów nietypowych;
nie można przekazać do niego dowolnej wartości z czasu wykonania.

\section{Przekazywanie doskonałe (\textit{perfect forwarding})}
\label{sec:variadic-forwarding}

Jednym z najważniejszych zastosowań pakietów jest budowanie
\emph{funkcji pośredniczących}, które przekazują argumenty do
innej funkcji lub konstruktora, możliwie wiernie zachowując
ich kategorię wartości (\emph{value category}).

\subsection{Różnica między lvalue a rvalue}

\begin{lstlisting}
#include <iostream>

void use(int&)  { std::cout << "lvalue\n"; }
void use(int&&) { std::cout << "rvalue\n"; }

int main() {
    int x = 10;
    use(x);       // lvalue
    use(20);      // rvalue
}
\end{lstlisting}

Parametr funkcji, choćby zadeklarowany jako \texttt{T\&\&}, jest
\emph{nazwanym obiektem}; użycie jego nazwy w wyrażeniu ma
kategorię \emph{lvalue}. Dlatego samo \texttt{use(arg)}
nie zachowa kategorii pierwotnego argumentu.

\subsection{Referencje przekazujące i \texttt{std::forward}}

\begin{lstlisting}
#include <utility>

template<class T>
void wrapper(T&& arg) {
    use(std::forward<T>(arg));
}

void forwarding_example() {
    int y = 0;
    wrapper(y);  // T = int&, forwarded as lvalue
    wrapper(30); // T = int, forwarded as rvalue
}
\end{lstlisting}

\texttt{T\&\&} jest \emph{referencją przekazującą}
(ang. \emph{forwarding reference}), gdy \texttt{T} jest
odpowiednio dedukowanym, niekwalifikowanym parametrem typu.
Działa tu \emph{składanie referencji}:

\[
T\&\;\& \longrightarrow T\&,
\qquad
T\&\;\&\& \longrightarrow T\&,
\qquad
T\&\&\;\&\& \longrightarrow T\&\&.
\]

W zapisie \texttt{std::forward<T>(arg)} jawny argument
\texttt{T} jest konieczny, by odtworzyć pierwotną kategorię.
\texttt{std::forward} nie wykonuje bezwarunkowego przeniesienia:
przepuszcza lvalue jako lvalue, a rvalue jako rvalue.

\subsection{Pakiet referencji przekazujących}

\begin{lstlisting}
#include <functional>
#include <utility>

template<class F, class... Args>
decltype(auto) call(F&& f, Args&&... args) {
    return std::invoke(
        std::forward<F>(f),
        std::forward<Args>(args)...
    );
}
\end{lstlisting}

Tutaj rozwijany jest \emph{cały wzorzec}
\texttt{std::forward<Args>(args)}, a nie tylko nazwa
\texttt{args}. Każdy element używa właściwego typu z pakietu
\texttt{Args}.
\texttt{std::invoke} potrafi wywoływać zwykłe funkcje,
obiekty funkcyjne i odpowiednio obsłużone wskaźniki do składowych.
\texttt{decltype(auto)} zachowuje typ zwracany przez docelowe
wywołanie, w tym referencję, jeżeli docelowa funkcja ją zwraca.

\begin{lstlisting}
#include <string>

std::string join(std::string a, std::string b) {
    return a + b;
}

auto result = call(join, std::string{"ab"}, "cd");
// result == "abcd"
\end{lstlisting}

\subsection{Czego nie robić: \texttt{std::move} i nazwy parametrów}

\begin{lstlisting}
template<class T>
void wrong_1(T&& arg) {
    use(arg);              // arg is always an lvalue here
}

template<class T>
void wrong_2(T&& arg) {
    use(std::move(arg));   // forces an rvalue even for lvalues
}
\end{lstlisting}

\texttt{std::move} to rzutowanie umożliwiające potraktowanie
obiektu jak rvalue, a nie polecenie ``zachowaj kategorię''.
Stosuj je, gdy rzeczywiście przekazujesz obiekt do operacji,
która może pobrać z niego zasób. W uniwersalnym wrapperze
używaj \texttt{std::forward}.

\paragraph{Ograniczenia żywotności.}
Referencje przekazujące \emph{nie przedłużają automatycznie}
żywotności obiektu do czasu dowolnego późniejszego wykonania.
Przechowanie referencji do argumentu tymczasowego w strukturze,
która przeżyje wywołanie, może prowadzić do wskaźników lub
referencji wiszących. Ten sam problem dotyczy między innymi
\texttt{std::forward\_as\_tuple}.

\section{Tworzenie obiektów z dowolnych argumentów}
\label{sec:variadic-emplace}

\subsection{Własna fabryka obiektów}

\begin{lstlisting}
#include <memory>
#include <utility>

template<class T, class... Args>
std::unique_ptr<T> create(Args&&... args) {
    return std::make_unique<T>(
        std::forward<Args>(args)...
    );
}

struct Point {
    int x, y;
    Point(int a, int b) : x(a), y(b) {}
};

auto point = create<Point>(4, 7);
\end{lstlisting}

\texttt{std::make\_unique} już dostarcza taką funkcjonalność,
więc \texttt{create} ma tu charakter dydaktyczny. Pokazuje
jednak wzorzec używany w fabrykach oraz funkcjach
\texttt{emplace} i \texttt{make}.

\subsection{Tworzenie kontenera bez wymogu kopiowalności}

\begin{lstlisting}
#include <concepts>
#include <string>
#include <utility>
#include <vector>

template<class T, class... Args>
    requires (std::constructible_from<T, Args&&> && ...)
std::vector<T> make_vector(Args&&... args) {
    std::vector<T> result;
    result.reserve(sizeof...(Args));
    (result.emplace_back(std::forward<Args>(args)), ...);
    return result;
}

auto words = make_vector<std::string>("A", "BB", "CCC");
\end{lstlisting}

Operator przecinka daje tu jednoznaczną kolejność wywołań
\texttt{emplace\_back}. Funkcja przyjmuje także pusty pakiet
argumentów. W przeciwieństwie do budowania wektora z listy
inicjalizacyjnej pozwala łatwo przekazywać do kolejnych
\texttt{emplace\_back} argumenty przenoszalne, ale niekopiowalne.
Dla elementów konstruowanych z więcej niż jednego parametru
wymagałaby innego interfejsu grupowania argumentów.

\section{\texttt{std::tuple}, \texttt{std::apply} i \texttt{std::index\_sequence}}
\label{sec:variadic-tuple}

\subsection{Heterogeniczna krotka \texttt{std::tuple}}

\texttt{std::tuple<Ts...>} to kontener o strukturze opisanej
listą typów znaną w czasie kompilacji. Każda pozycja może
mieć inny typ:

\begin{lstlisting}
#include <tuple>
#include <string>

std::tuple<int, std::string, double> t{7, "Ada", 3.5};

int n = std::get<0>(t);
const std::string& name = std::get<1>(t);
\end{lstlisting}

W przeciwieństwie do \texttt{std::vector<T>}, krotka nie
przechowuje dowolnej liczby elementów jednego typu, lecz
\emph{ustaloną na etapie kompilacji} listę typów.

\subsection{Rozwinięcie krotki do argumentów funkcji}

Funkcja \texttt{std::apply} (C++17) wywołuje funkcję,
przekazując elementy obiektu \emph{tuple-like} jako osobne
argumenty:

\begin{lstlisting}
#include <iostream>
#include <tuple>

void apply_example() {
    std::apply([](const auto&... xs) {
        ((std::cout << xs << ' '), ...);
        std::cout << '\n';
    }, t);
}
\end{lstlisting}

\texttt{std::apply} oszczędza ręcznego konstruowania
sekwencji indeksów, ale warto wiedzieć, jak taki mechanizm
można zaimplementować.

\subsection{Pakiety indeksów: \texttt{std::index\_sequence}}

\begin{lstlisting}
#include <cstddef>
#include <iostream>
#include <tuple>
#include <utility>

template<class Tuple, std::size_t... I>
void print_tuple_impl(const Tuple& t,
                      std::index_sequence<I...>) {
    ((std::cout << std::get<I>(t) << ' '), ...);
    std::cout << '\n';
}

template<class... Ts>
void print_tuple(const std::tuple<Ts...>& t) {
    print_tuple_impl(t, std::index_sequence_for<Ts...>{});
}

void tuple_example() {
    print_tuple(std::make_tuple(1, "abc", 3.0));
}
\end{lstlisting}

\texttt{std::index\_sequence<0,1,2>} jest specjalizacją
szablonu przyjmującego pakiet indeksów całkowitych.
\texttt{std::index\_sequence\_for<Ts...>} generuje indeksy
$0,1,\ldots,n-1$ dla pakietu długości $n$.
Wzorzec \texttt{std::get<I>(t)...} rozwija wyrażenie
osobno dla każdego indeksu.

\subsection{Równoległe rozwijanie pakietów}

W jednym wzorcu można użyć kilku pakietów. Muszą mieć
\emph{tę samą długość}, jeśli są rozwijane równolegle:

\begin{lstlisting}
#include <tuple>
#include <type_traits>
#include <utility>

template<class A, class B>
struct zipped_types;

template<class... Ts, class... Us>
struct zipped_types<std::tuple<Ts...>, std::tuple<Us...>> {
    static_assert(sizeof...(Ts) == sizeof...(Us));
    using type = std::tuple<std::pair<Ts, Us>...>;
};

using Z = zipped_types<
    std::tuple<int, char>,
    std::tuple<double, bool>
>::type;

static_assert(std::is_same_v<
    Z, std::tuple<std::pair<int, double>,
                  std::pair<char, bool>>>);
\end{lstlisting}

Elementy \texttt{Ts} i \texttt{Us} są łączone pozycjami:
$\bigl(T_1,U_1\bigr),\ldots,\bigl(T_n,U_n\bigr)$.
Nie jest to iloczyn kartezjański pakietów.

\section{Dziedziczenie wariadyczne i \texttt{std::visit}}
\label{sec:variadic-overloaded}

Pakiet \texttt{Bases...} można rozwinąć na liście baz.
Z kolei zapis \texttt{using Bases::operator()...}
wprowadza ich operatory wywołania do jednej klasy.
W ten sposób otrzymujemy zestaw przeciążeń.

\begin{lstlisting}
#include <iostream>
#include <string>
#include <variant>

template<class... Fs>
struct overloaded : Fs... {
    using Fs::operator()...;
};

template<class... Fs>
overloaded(Fs...) -> overloaded<Fs...>; // C++17 deduction guide

int main() {
    std::variant<int, std::string> value = "hello";

    std::visit(overloaded{
        [](int n) { std::cout << "int: " << n << '\n'; },
        [](const std::string& s) {
            std::cout << "string: " << s << '\n';
        }
    }, value);
}
\end{lstlisting}

Jest to standardowy idiom C++17: każda lambda definiuje
osobny \texttt{operator()}, a dziedziczenie wariadyczne
łączy je w jeden obiekt wywoływalny. Gdy stosujemy ten wzorzec
w wielu miejscach, warto umieścić \texttt{overloaded} w
niewielkim nagłówku narzędziowym.

\section{Ograniczanie pakietów za pomocą \texttt{concepts} (C++20)}
\label{sec:variadic-concepts}

Nie każda operacja ma sens dla każdej możliwej kombinacji
typów. W C++20 można czytelnie określać wymagania dla
\emph{wszystkich} elementów pakietu.

\begin{lstlisting}
#include <concepts>
#include <type_traits>

template<class... Ts>
    requires (std::integral<Ts> && ...)
constexpr bool all_nonnegative(Ts... xs) {
    return ((xs >= 0) && ...);
}

static_assert(all_nonnegative(0, 2, 7));
static_assert(!all_nonnegative(0, -1, 7));

// all_nonnegative(3.14); // rejected by constraints
\end{lstlisting}

Możemy również ograniczać pakiet w odniesieniu do
konkretnego typu docelowego:

\begin{lstlisting}
#include <concepts>

template<class T, class... Args>
concept ConstructibleWith =
    std::constructible_from<T, Args...>;

static_assert(ConstructibleWith<std::string, const char*>);
static_assert(!ConstructibleWith<int, std::string>);
\end{lstlisting}

Pierwszy przykład stosuje fold nad \emph{warunkami dla pojedynczych
typów}; drugi sprawdza \emph{jedną własność całej sekwencji
typów}, tj. czy odpowiedni konstruktor przyjmuje argumenty
\texttt{Args...} w tej kolejności. Tych dwóch mechanizmów
nie należy utożsamiać.

\subsection{\texttt{if constexpr} i struktura pakietu}

\begin{lstlisting}
#include <iostream>

template<class... Ts>
void report(const Ts&... xs) {
    if constexpr (sizeof...(Ts) == 0) {
        std::cout << "empty\n";
    } else {
        ((std::cout << xs << ' '), ...);
        std::cout << '\n';
    }
}
\end{lstlisting}

\texttt{if constexpr} pozwala w zależności od długości pakietu
wybierać gałąź w czasie kompilacji. W dobrze sformułowanym
szablonie odrzucona gałąź nie jest instancjonowana dla danej
specjalizacji. To użyteczne przy przypadkach brzegowych,
które nie mają sensu dla pustego pakietu.

\section{Ograniczenia, pułapki i diagnostyka}
\label{sec:variadic-pulapki}

\subsection{Pakiet typów nie jest pakietem argumentów funkcji}

\begin{lstlisting}
template<class... Ts>
void f(Ts... args) {
    // sizeof...(Ts): number of types
    // sizeof...(args): number of function parameters
    // args[0];  // ERROR: args is not an array
}
\end{lstlisting}

Do wybierania argumentu na ustalonej pozycji można użyć
\texttt{std::tuple} i \texttt{std::get} albo wywoływać funkcję
przez \texttt{std::apply}. Sam pakiet \texttt{args...}
nie jest indeksowanym kontenerem.

\subsection{Brak zgodności typów po rozwinięciu}

\begin{lstlisting}
template<class... Ts>
auto add(Ts... xs) {
    return (0 + ... + xs);
}

// add(1, std::string{"x"}); // ERROR: unsuitable operator+
\end{lstlisting}

Pakiet nie omija systemu typów. Każde wyrażenie powstałe po
jego rozwinięciu musi być poprawne, podobnie jak zwykły kod
C++. Jeżeli szablon ma określone wymagania, najlepiej
zapisać je za pomocą \texttt{requires}.

\subsection{Niewłaściwe przekazywanie kategorii wartości}

\begin{lstlisting}
template<class F, class... Args>
void bad_call(F&& f, Args&&... args) {
    std::forward<F>(f)(args...); // loses rvalue categories
}

// Correct:
// std::forward<F>(f)(std::forward<Args>(args)...);
\end{lstlisting}

Dla typów niekopiowalnych taki błąd często ujawnia się
natychmiast podczas kompilacji. Dla typów kopiowalnych
może spowodować jedynie dodatkowe kopie albo zmianę
wybranego przeciążenia.

\subsection{Wnioskowanie typów i \texttt{\{...\}}}

\begin{lstlisting}
#include <initializer_list>

template<class... Args>
void foo(Args&&...) {}

void braced_example() {
    // foo({1, 2, 3}); // ERROR: cannot deduce Args
    foo(std::initializer_list<int>{1, 2, 3}); // OK
}
\end{lstlisting}

Sama lista ujęta w nawiasy klamrowe \texttt{\{...\}} nie ma
typu, który dowolny pakiet \texttt{Args\&\&...} potrafiłby
wydedukować. Potrzebny jest określony typ parametru,
np. \texttt{std::initializer\_list}\allowbreak\texttt{<int>}, lub jawnie
utworzony obiekt.

\subsection{Ograniczenia położenia pakietu}

W \emph{głównej definicji szablonu klasy} pakiet parametrów
szablonu musi być ostatni. Dla szablonów funkcji istnieją
bardziej elastyczne reguły, ale kolejne parametry muszą
być dedukowalne lub mieć argumenty domyślne, stosownie do
kontekstu:

\begin{lstlisting}
// ERROR for a primary class template:
// template<class... Ts, class U> struct invalid;

template<class U, class... Ts>
struct valid {};
\end{lstlisting}

\subsection{Skomplikowany kod zamiast prostego rozwiązania}

Szablon wariadyczny to narzędzie, a nie obowiązkowy styl.
Jeżeli argumenty mają \emph{jednolity typ i ich liczba
zmienia się w czasie wykonania}, lepszy bywa
\texttt{std::span<T>}, \texttt{std::vector<T>} lub para
iteratorów. Pakiet ma liczbę elementów i strukturę typów
ustalaną na etapie kompilacji; nie zastępuje kontenera.

\section{Variadic templates a funkcje wariadyczne w stylu C}

Historyczne funkcje z deklaracją \texttt{void f(int, ...)}
korzystają z mechanizmu argumentów wariadycznych języka C,
a nie z pakietu parametrów szablonu C++:

\begin{lstlisting}
#include <cstdio>

int main() {
    std::printf("%d %f\n", 7, 2.5); // C-style variadic call
}
\end{lstlisting}

W formacie C zgodność argumentów ze specyfikatorami
formatowania jest obowiązkiem programisty (niektóre
kompilatory potrafią wykrywać część pomyłek). W szablonach
wariadycznych typy argumentów pozostają dostępne
kompilatorowi, dzięki czemu można wymuszać ograniczenia
i stosować przeciążanie oraz \texttt{std::forward}.
Nie należy jednak wyciągać z tego wniosku, że dowolny
szablon automatycznie bezbłędnie formatuje argumenty.

\section{Pełny przykład: prosty rejestrator zdarzeń}
\label{sec:variadic-projekt}

Zastosujmy omówione mechanizmy do funkcji tworzącej tekstowy
komunikat z dowolnej liczby argumentów możliwych do
wypisania do strumienia. Zamiast ręcznie doklejać spacje
w wielu przeciążeniach, użyjemy fold nad operatorem przecinka.

\begin{lstlisting}
#include <iostream>
#include <sstream>
#include <string>

// Each argument must support insertion into std::ostream.
template<class... Args>
std::string make_message(const Args&... args) {
    std::ostringstream out;
    ((out << args), ...);
    return out.str();
}

template<class... Args>
void log_line(const Args&... args) {
    std::cout << "[LOG] " << make_message(args...) << '\n';
}

int main() {
    log_line("count=", 3, ", ratio=", 0.75);
    log_line("finished");
    log_line();
}
\end{lstlisting}

\paragraph{Analiza.}
Dla argumentów \texttt{("count=",3,", ratio=",0.75)}
pakiet zawiera cztery elementy. Wyrażenie
\texttt{((out << args), ...)} wstawia je kolejno do strumienia,
a \texttt{make\_message} zwraca powstały napis. Dla
pustego pakietu fold z operatorem przecinka jest poprawny,
a wynikowy napis jest pusty. W wersji produkcyjnej należałoby
rozważyć wymagania dotyczące formatowania, zachowanie w
obliczu wyjątków, współbieżność oraz ewentualne skorzystanie
z narzędzia formatującego wspieranego przez używany kompilator.

\section{Zestawienie: kiedy używać poszczególnych mechanizmów?}

\begin{center}
\small
\begin{tabular}{|p{0.40\textwidth}|p{0.46\textwidth}|}
\hline
\textbf{Cel} & \textbf{Typowe rozwiązanie} \\
\hline
Dowolna liczba argumentów o różnych typach & \texttt{template<class... Ts>} i \texttt{Ts...} \\
\hline
Wykonanie tej samej operacji dla każdego argumentu & Fold z operatorem \texttt{,} (C++17) \\
\hline
Suma, iloczyn lub wspólny warunek & Fold z \texttt{+}, \texttt{*}, \texttt{\&\&} lub \texttt{||} \\
\hline
Przekazanie argumentów dalej z zachowaniem kategorii & \texttt{Args\&\&...}, \texttt{std::forward<Args>(args)...} \\
\hline
Zapis heterogenicznych wartości & \texttt{std::tuple<Ts...>} \\
\hline
Wywołanie funkcji elementami krotki & \texttt{std::apply} (C++17) \\
\hline
Rozwinięcie krotki przy pomocy indeksów & \texttt{std::index\_sequence} \\
\hline
Lista typów do obliczeń podczas kompilacji & Wariadyczny szablon klasy i specjalizacje \\
\hline
Ograniczenie do odpowiednich typów & Fold z koncepcjami i \texttt{requires} (C++20) \\
\hline
Zmienna długość listy jednolitych danych w czasie wykonania & \texttt{std::vector} lub \texttt{std::span}, nie pakiet \\
\hline
\end{tabular}
\end{center}

\section{Podsumowanie i zadania do samodzielnego sprawdzenia}

\begin{enumerate}
    \item \texttt{class... Ts} deklaruje \emph{pakiet typów},
    a \texttt{Ts... xs} wprowadza \emph{pakiet parametrów funkcji}.
    \item \texttt{sizeof...(Ts)} oblicza liczbę elementów pakietu
    podczas kompilacji.
    \item \texttt{pattern...} rozwija cały \emph{wzorzec},
    nie tylko pojedynczy identyfikator.
    \item Fold expressions (od C++17) zastępują wiele prostych
    rekurencyjnych implementacji z C++11.
    \item Dla pustego pakietu istotne są element neutralny i
    reguły operatora; nie każda forma fold jest poprawna.
    \item W wrapperach zachowuj kategorię wartości argumentów,
    używając referencji przekazujących i \texttt{std::forward}.
    \item \texttt{std::tuple}, \texttt{std::apply} i
    \texttt{std::index\_sequence} współpracują z pakietami,
    umożliwiając operacje na heterogenicznych danych.
    \item W C++20 używaj \texttt{concepts}/\texttt{requires},
    aby błędne wywołania odrzucać na granicy interfejsu.
\end{enumerate}

\paragraph{Pytania kontrolne.}
Spróbuj przewidzieć wyniki poniższych wyrażeń bez uruchamiania
programu. Następnie zweryfikuj je za pomocą \texttt{static\_assert}:

\begin{lstlisting}
// static_assert(sizeof...(int) == 1); // ERROR: int is not a pack

// For the functions defined above:
static_assert(product() == 1);
static_assert(left_sub(10, 3) == 87);
static_assert(right_sub(10, 3) == 107);
static_assert(count_integral_v<int, double, char> == 2);
\end{lstlisting}

Pierwsza linia jest \emph{celowo niepoprawna}: operator
\texttt{sizeof...} wymaga nazwy pakietu, a nie zwykłego typu.

% Wykorzystane punkty odniesienia (bez kopiowania tresci):
% https://en.cppreference.com/w/cpp/language/parameter_pack
% https://en.cppreference.com/w/cpp/language/fold
% https://en.cppreference.com/w/cpp/utility/forward
% https://en.cppreference.com/w/cpp/utility/apply
% https://en.cppreference.com/w/cpp/utility/integer_sequence


\chapter[Variadic templates w C++]{Szablony o zmiennej liczbie argumentów (\textit{variadic templates})}
\label{chap:variadic-templates-cpp}

\section{Idea i zastosowania}

Zwykły szablon funkcji przyjmuje z góry ustaloną liczbę parametrów
szablonu. \emph{Variadic template} (szablon o zmiennej liczbie parametrów)
może natomiast przyjmować \emph{dowolną liczbę} argumentów szablonu,
również zero. Mechanizmem, który to umożliwia, jest
\emph{pakiet parametrów} (ang. \emph{parameter pack}).

\begin{lstlisting}
template<class T>
void single(const T&) {}        // one type

template<class... Ts>
void many(const Ts&...) {}      // zero or more types

void demo() {
    many();                     // zero arguments
    many(10);                   // one argument
    many(10, 3.5, "abc");      // three arguments
}
\end{lstlisting}

Pakiety parametrów są ważną częścią nowoczesnego C++.
Wykorzystują je m.in. klasy \texttt{std::tuple} i \texttt{std::variant}
oraz fabryki \texttt{std::make\_unique} i \texttt{std::make\_shared}.
Nie trzeba dzięki nim pisać przeciążeń dla każdej
liczby argumentów.

\paragraph{Wymagane wersje języka.}
Pakiety parametrów i perfect forwarding pojawiły się w C++11.
\emph{Fold expressions} dodano w C++17, a koncepcje
(\texttt{concepts}) i klauzulę \texttt{requires} w C++20.
Dalej zakładamy C++20.

\section{Trzy kropki: deklaracja pakietu a rozwinięcie}

\begin{center}\small
\begin{tabularx}{\textwidth}{@{}p{0.35\textwidth}X@{}}
\toprule
Zapis & Działanie \\
\midrule
\code{typename... Ts} & Deklaruje pakiet parametrów typowych szablonu. \\
\code{Ts...} & Rozwija pakiet typów w miejscu, które dopuszcza taki wzorzec. \\
\code{args...} & Rozwija pakiet argumentów, tworząc listę wyrażeń. \\
\code{sizeof...(Ts)} & Podaje liczbę elementów pakietu w czasie kompilacji. \\
\code{std::forward<Ts>(args)...} & Przekazuje każdy argument z zachowaniem jego kategorii wartości. \\
\bottomrule
\end{tabularx}
\end{center}
\label{sec:variadic-trzy-kropki}

W różnych miejscach token \texttt{...} oznacza różne operacje.
Rozróżnienie to stanowi podstawę poprawnego używania szablonów
wariadycznych.

\begin{lstlisting}
template<class... Ts>             // declare type pack Ts
void f(Ts... xs) {                // declare function parameter pack xs
    // ...
}
\end{lstlisting}

Deklaracja \texttt{class... Ts} wprowadza pakiet \emph{typów}.
Zapis \texttt{Ts... xs} rozwija pakiet typów i wprowadza odpowiadający
mu pakiet parametrów funkcji. Wywołanie
\texttt{f(1, 2.0, 'x')} tworzy w uproszczeniu funkcję o parametrach
\texttt{(int, double, char)}.

\subsection{Rozwinięcie pakietu (\textit{pack expansion})}

Rozwinięcie składa się ze \emph{wzorca} (\emph{pattern}) oraz następującego
po nim \texttt{...}. Dla każdego elementu pakietu kompilator podstawia
kolejny element do całego wzorca:

\begin{lstlisting}
template<class... Ts>
void demo(Ts... xs) {
    g(xs...);                // g(x1, x2, ..., xn)
    h((xs * 2)...);          // h(x1*2, x2*2, ..., xn*2)
}
\end{lstlisting}

W drugim wywołaniu wzorcem jest \texttt{(xs * 2)}, a \emph{nie}
pojedyncza zmienna \texttt{xs}. Ważne: \texttt{xs...} nie jest
kontenerem ani tablicą; nie można po nim bezpośrednio iterować zwykłą
pętlą \texttt{for}. To składnia kompilatora służąca do generowania
listy wyrażeń.

\paragraph{Pakiet typów a pakiet wartości.}
W przykładzie \texttt{Ts} reprezentuje typy, natomiast
\texttt{xs} to konkretne parametry funkcji. Jeden typ może wystąpić
w pakiecie wielokrotnie, a każdy element pakietu wartości zachowuje
swój właściwy typ.

\subsection{Liczba elementów: \texttt{sizeof...}}

Operator \texttt{sizeof...} daje liczbę elementów pakietu jako
stałą czasu kompilacji. Nie należy go mylić z operatorem
\texttt{sizeof} mierzącym rozmiar typu lub obiektu.

\begin{lstlisting}
#include <cstddef>

template<class... Ts>
constexpr std::size_t count(Ts&&... xs) {
    static_assert(sizeof...(Ts) == sizeof...(xs));
    return sizeof...(xs);
}

static_assert(count() == 0);
static_assert(count(2, 3.5, 'a') == 3);
\end{lstlisting}

Nazwy \texttt{xs} nie używamy w ciele funkcji do odczytywania
wartości; operacja liczenia jest czysto składniowa.

\section{Pierwsze praktyczne funkcje}

\subsection{Wypisywanie dowolnej liczby argumentów}

Najprostszy sposób wykorzystania pakietu to jego rozwinięcie
w wywołaniu innej funkcji:

\begin{lstlisting}
#include <iostream>
#include <string>

void output(int x, double y, const std::string& z) {
    std::cout << x << ' ' << y << ' ' << z << '\n';
}

template<class... Args>
void forward_to_output(Args... args) {
    output(args...);
}

int main() {
    forward_to_output(3, 2.5, std::string{"hello"});
}
\end{lstlisting}

Funkcja jest poprawna tylko dla takich wywołań, z których można
utworzyć poprawne wywołanie \texttt{output}. Sama wariadyczność
nie oznacza, że każda liczba i każdy rodzaj argumentów będą
akceptowane przez funkcję docelową.

\subsection{Przekazywanie przez wartość i przez referencję}

\begin{lstlisting}
template<class... Ts>
void by_value(Ts... xs);            // may copy/move

template<class... Ts>
void by_const_ref(const Ts&... xs); // read-only; no copies

template<class... Ts>
void by_ref(Ts&... xs);             // needs lvalues
\end{lstlisting}

Wybór postaci parametrów ma te same skutki co w zwykłych funkcjach.
W szczególności \texttt{Ts... xs} tworzy osobne parametry przez wartość,
a \texttt{const Ts\&... xs} zwykle pozwala uniknąć niepotrzebnego
kopiowania. Przypadek \texttt{Ts\&\&... xs} omówimy dokładniej
w sekcji o \emph{perfect forwarding}.

\section{Rekurencyjne przetwarzanie pakietu: metoda C++11}
\label{sec:variadic-rekurencja}

Przed C++17 często rozbijano pakiet na \emph{pierwszy argument}
i \emph{pozostałe argumenty}, a następnie wywoływano funkcję
rekurencyjnie. Potrzebny był osobny przypadek kończący rekurencję.

\begin{lstlisting}
#include <iostream>

void print_recursive() {
    std::cout << '\n';        // base case
}

template<class First, class... Rest>
void print_recursive(const First& first,
                     const Rest&... rest) {
    std::cout << first << ' ';
    print_recursive(rest...); // one argument fewer
}

int main() {
    print_recursive(1, "two", 3.0);
}
\end{lstlisting}

Przy trzech argumentach instancjonowane są kolejno wersje
z pakietami o długościach $3$, $2$, $1$ oraz przypadek końcowy.
Jest to \emph{rekurencja szablonowa}; nie należy automatycznie
utożsamiać jej z dynamiczną rekurencją wykonywaną w programie.
Kompilator może wygenerować zwykłą sekwencję operacji.

\paragraph{Kiedy stosować?}
Rekurencyjne rozkładanie pakietu nadal bywa przydatne, jeśli kolejne
kroki mają odrębne warunki lub należy budować typy rekurencyjne.
Jeżeli wykonujemy tę samą operację dla wszystkich argumentów,
w C++17 i nowszych zazwyczaj lepszym wyborem jest \emph{fold expression}.

\section{Fold expressions (C++17)}

\begin{center}\small
\begin{tabularx}{\textwidth}{@{}p{0.35\textwidth}X@{}}
\toprule
Forma złożenia & Sposób rozwinięcia \\
\midrule
\code{(... + args)} & Złożenie lewostronne: kolejne dodawania łączone od lewej. \\
\code{(args + ...)} & Złożenie prawostronne: kolejne dodawania łączone od prawej. \\
\code{(init + ... + args)} & Złożenie z wartością początkową; działa także dla pustego pakietu. \\
\code{(f(args), ...)} & Wywołuje \code{f} dla każdego argumentu od lewej do prawej dzięki operatorowi przecinka. \\
\texttt{(args \&\& ...)} & Łączy warunki z krótkim spięciem; dla pustego pakietu daje \code{true}. \\
\bottomrule
\end{tabularx}
\end{center}
\label{sec:variadic-fold}

\emph{Fold expression} redukuje pakiet wyrażeń z użyciem jednego
operatora binarnego, np. \texttt{+}, \texttt{\&\&}, \texttt{||}
lub operatora przecinka. Nie wymaga jawnego przypadku
rekurencyjnego.

\subsection{Cztery postacie składni}

Niech \texttt{xs} oznacza pakiet, $\circ$ ustalony operator,
a $z$ wartość początkową. Istnieją cztery postacie składni:

\begin{center}
\small
\begin{tabular}{|p{0.39\textwidth}|p{0.47\textwidth}|}
\hline
\textbf{Postać C++} & \textbf{Sposób łączenia dla trzech argumentów} \\
\hline
\texttt{(xs op ...)} & $x_1\circ(x_2\circ x_3)$ (prawostronne, bez inicjalizatora) \\
\hline
\texttt{(... op xs)} & $(x_1\circ x_2)\circ x_3$ (lewostronne, bez inicjalizatora) \\
\hline
\texttt{(xs op ... op z)} & $x_1\circ(x_2\circ(x_3\circ z))$ (prawostronne, z inicjalizatorem) \\
\hline
\texttt{(z op ... op xs)} & $((z\circ x_1)\circ x_2)\circ x_3$ (lewostronne, z inicjalizatorem) \\
\hline
\end{tabular}
\end{center}

W tabeli \texttt{op} jest jedynie oznaczeniem miejsca na rzeczywisty
operator, np. \texttt{+}; nie jest słowem kluczowym języka.
\emph{Nawiasy całego fold expression są obowiązkowe}.
Dwa operatory występujące w fold z inicjalizatorem muszą być takie same.

\subsection{Suma, iloczyn i warunki logiczne}

\begin{lstlisting}
#include <concepts>

template<class... Ts>
    requires (std::integral<Ts> && ...)
constexpr auto sum(Ts... xs) {
    return (0 + ... + xs);       // left fold with identity 0
}

template<class... Ts>
constexpr bool all_true(Ts... xs) {
    return (true && ... && static_cast<bool>(xs));
}

template<class... Ts>
constexpr bool any_true(Ts... xs) {
    return (false || ... || static_cast<bool>(xs));
}

static_assert(sum() == 0);
static_assert(sum(2, 3, 4) == 9);
static_assert(all_true());
static_assert(!any_true());
static_assert(all_true(1, 2, 3));
\end{lstlisting}

Wyrażenie \texttt{(0 + ... + xs)} może być używane także dla
pustego pakietu. Początkowe \texttt{0} wpływa jednak na typ wyniku
i ewentualne promocje arytmetyczne; nie jest to uniwersalna
implementacja sumy dowolnych typów liczbowych.

\subsection{Lewostronne i prawostronne składanie nie zawsze są równe}

Dla operatora odejmowania kolejność nawiasowania zmienia wynik:

\begin{lstlisting}
template<class... Ts>
constexpr auto left_sub(Ts... xs) {
    return (100 - ... - xs);
}

template<class... Ts>
constexpr auto right_sub(Ts... xs) {
    return (xs - ... - 100);
}

static_assert(left_sub(10, 3) == 87);  // (100 - 10) - 3
static_assert(right_sub(10, 3) == 107); // 10 - (3 - 100)
\end{lstlisting}

Fold określa \emph{strukturę nawiasów}, ale nie należy mylić tej
własności z ogólną gwarancją kolejności wykonywania wszystkich
podwyrażeń dla dowolnego operatora. Reguły ewaluacji wynikają
również z semantyki danego operatora.

\subsection{Pusty pakiet i element neutralny}

Dla \emph{jednoargumentowego} fold z pustym pakietem C++ dopuszcza
zaledwie trzy operatory: \texttt{\&\&} zwraca \texttt{true},
\texttt{||} zwraca \texttt{false}, a \texttt{,} daje
wyrażenie typu \texttt{void}. Pusty pakiet w
\texttt{(... + xs)} nie jest poprawny. Stosując wersję
z wartością początkową, możemy jawnie obsłużyć pusty pakiet:

\begin{lstlisting}
template<class... Ts>
constexpr auto product(Ts... xs) {
    return (1 * ... * xs);
}

static_assert(product() == 1);
static_assert(product(2, 3, 4) == 24);
\end{lstlisting}

\section{Operator przecinka i kolejność wykonywania}
\label{sec:variadic-przecinek}

Operator przecinka jest wygodny, gdy dla każdego elementu pakietu
chcemy wykonać instrukcję o efektach ubocznych. W odróżnieniu
od kolejności ewaluacji zwykłych argumentów funkcji, jego
wbudowana postać gwarantuje wykonanie lewej strony przed prawą.

\begin{lstlisting}
#include <iostream>

template<class... Ts>
void print_all(const Ts&... xs) {
    ((std::cout << xs << ' '), ...);
    std::cout << '\n';
}

int main() {
    print_all(10, 2.5, "hello");
    print_all();                // only a newline
}
\end{lstlisting}

Zewnętrzne nawiasy tworzą fold, a wewnętrzne oddzielają
wyrażenie wypisujące od operatora \texttt{,}.

\paragraph{Kolejność argumentów zwykłej funkcji.}

\begin{lstlisting}
template<class... Ts>
void sink(Ts...);

// sink(change_a(), change_b(), change_c());
\end{lstlisting}

Samo rozwinięcie \texttt{...} w liście argumentów funkcji
\emph{nie} gwarantuje, że \texttt{change\_a()} zostanie obliczone
przed \texttt{change\_b()}. W C++17 i nowszych wywołania
te nie przeplatają swoich efektów, ale ich wzajemna kolejność
nie jest określona. Gdy kolejność jest istotna, można użyć
operatora przecinka lub innej konstrukcji dającej odpowiednie
gwarancje sekwencjonowania.

\section{Pakiety typów i metaprogramowanie}
\label{sec:variadic-type-list}

Wariadyczność może dotyczyć nie tylko funkcji, lecz także
\emph{szablonów klas}. W takim przypadku pakiet przechowuje
listę typów, a instancjonowanie klasy odbywa się podczas kompilacji.

\subsection{Własna lista typów}

\begin{lstlisting}
#include <cstddef>
#include <type_traits>

template<class... Ts>
struct type_list {};

template<class List>
struct list_size;

template<class... Ts>
struct list_size<type_list<Ts...>>
    : std::integral_constant<std::size_t, sizeof...(Ts)> {};

using Types = type_list<int, double, char>;
static_assert(list_size<Types>::value == 3);
static_assert(list_size<type_list<>>::value == 0);
\end{lstlisting}

Jest to przykład \emph{częściowej specjalizacji} szablonu klasy.
Wzorzec \texttt{type\_list<Ts...>} dopasowuje dowolną listę
typów, a \texttt{sizeof...(Ts)} oblicza jej długość.

\subsection{Wydobywanie pierwszego typu}

\begin{lstlisting}
template<class List>
struct first_type;   // intentionally undefined for empty list

template<class T, class... Rest>
struct first_type<type_list<T, Rest...>> {
    using type = T;
};

using X = first_type<type_list<long, float>>::type;
static_assert(std::is_same_v<X, long>);
\end{lstlisting}

Rozkład na \texttt{T} oraz \texttt{Rest...} jest odpowiednikiem
rekurencyjnego rozdzielenia argumentów funkcji na pierwszy
i pozostałe. Dla pustej listy \texttt{type\_list<>} nazwa
\texttt{first\_type<...>::type} celowo nie istnieje.

\subsection{Czy wszystkie typy spełniają własność?}

\begin{lstlisting}
#include <type_traits>

template<class... Ts>
inline constexpr bool all_integral_v =
    (std::is_integral_v<Ts> && ...);

template<class... Ts>
inline constexpr std::size_t count_integral_v =
    (std::size_t{0} + ... +
        (std::is_integral_v<Ts> ? 1u : 0u));

static_assert(all_integral_v<int, long, char>);
static_assert(!all_integral_v<int, double>);
static_assert(count_integral_v<int, double, char> == 2);
\end{lstlisting}

Jest to częsty wzorzec: pakiet typów łączymy z cechami
typów (\emph{type traits}) i operatorami logicznymi fold.

\section{Pakiety wartości i parametry nietypowe}

Pakiet może także zawierać wartości znane w czasie kompilacji,
a nie typy. W szablonie klasy służą do tego
\emph{nietypowe parametry szablonu} (\emph{non-type template parameters}).

\begin{lstlisting}
#include <array>
#include <cstddef>

template<int... Ns>
struct int_pack {
    static constexpr std::size_t size = sizeof...(Ns);
    static constexpr int sum = (0 + ... + Ns);
};

static_assert(int_pack<>::sum == 0);
static_assert(int_pack<2, 4, 6>::sum == 12);
\end{lstlisting}

Od C++17 można zadeklarować pakiet \texttt{auto...}:

\begin{lstlisting}
template<auto... Values>
struct value_pack {
    static constexpr std::size_t size = sizeof...(Values);
};

static_assert(value_pack<42, 'x', true>::size == 3);
\end{lstlisting}

Dla \texttt{auto...} dopuszczalność konkretnej wartości jako
argumentu szablonu nadal podlega regułom parametrów nietypowych;
nie można przekazać do niego dowolnej wartości z czasu wykonania.

\section{Przekazywanie doskonałe (\textit{perfect forwarding})}
\label{sec:variadic-forwarding}

Jednym z najważniejszych zastosowań pakietów jest budowanie
\emph{funkcji pośredniczących}, które przekazują argumenty do
innej funkcji lub konstruktora, możliwie wiernie zachowując
ich kategorię wartości (\emph{value category}).

\subsection{Różnica między lvalue a rvalue}

\begin{lstlisting}
#include <iostream>

void use(int&)  { std::cout << "lvalue\n"; }
void use(int&&) { std::cout << "rvalue\n"; }

int main() {
    int x = 10;
    use(x);       // lvalue
    use(20);      // rvalue
}
\end{lstlisting}

Parametr funkcji, choćby zadeklarowany jako \texttt{T\&\&}, jest
\emph{nazwanym obiektem}; użycie jego nazwy w wyrażeniu ma
kategorię \emph{lvalue}. Dlatego samo \texttt{use(arg)}
nie zachowa kategorii pierwotnego argumentu.

\subsection{Referencje przekazujące i \texttt{std::forward}}

\begin{lstlisting}
#include <utility>

template<class T>
void wrapper(T&& arg) {
    use(std::forward<T>(arg));
}

void forwarding_example() {
    int y = 0;
    wrapper(y);  // T = int&, forwarded as lvalue
    wrapper(30); // T = int, forwarded as rvalue
}
\end{lstlisting}

\texttt{T\&\&} jest \emph{referencją przekazującą}
(ang. \emph{forwarding reference}), gdy \texttt{T} jest
odpowiednio dedukowanym, niekwalifikowanym parametrem typu.
Działa tu \emph{składanie referencji}:

\[
T\&\;\& \longrightarrow T\&,
\qquad
T\&\;\&\& \longrightarrow T\&,
\qquad
T\&\&\;\&\& \longrightarrow T\&\&.
\]

W zapisie \texttt{std::forward<T>(arg)} jawny argument
\texttt{T} jest konieczny, by odtworzyć pierwotną kategorię.
\texttt{std::forward} nie wykonuje bezwarunkowego przeniesienia:
przepuszcza lvalue jako lvalue, a rvalue jako rvalue.

\subsection{Pakiet referencji przekazujących}

\begin{lstlisting}
#include <functional>
#include <utility>

template<class F, class... Args>
decltype(auto) call(F&& f, Args&&... args) {
    return std::invoke(
        std::forward<F>(f),
        std::forward<Args>(args)...
    );
}
\end{lstlisting}

Tutaj rozwijany jest \emph{cały wzorzec}
\texttt{std::forward<Args>(args)}, a nie tylko nazwa
\texttt{args}. Każdy element używa właściwego typu z pakietu
\texttt{Args}.
\texttt{std::invoke} potrafi wywoływać zwykłe funkcje,
obiekty funkcyjne i odpowiednio obsłużone wskaźniki do składowych.
\texttt{decltype(auto)} zachowuje typ zwracany przez docelowe
wywołanie, w tym referencję, jeżeli docelowa funkcja ją zwraca.

\begin{lstlisting}
#include <string>

std::string join(std::string a, std::string b) {
    return a + b;
}

auto result = call(join, std::string{"ab"}, "cd");
// result == "abcd"
\end{lstlisting}

\subsection{Czego nie robić: \texttt{std::move} i nazwy parametrów}

\begin{lstlisting}
template<class T>
void wrong_1(T&& arg) {
    use(arg);              // arg is always an lvalue here
}

template<class T>
void wrong_2(T&& arg) {
    use(std::move(arg));   // forces an rvalue even for lvalues
}
\end{lstlisting}

\texttt{std::move} to rzutowanie umożliwiające potraktowanie
obiektu jak rvalue, a nie polecenie ``zachowaj kategorię''.
Stosuj je, gdy rzeczywiście przekazujesz obiekt do operacji,
która może pobrać z niego zasób. W uniwersalnym wrapperze
używaj \texttt{std::forward}.

\paragraph{Ograniczenia żywotności.}
Referencje przekazujące \emph{nie przedłużają automatycznie}
żywotności obiektu do czasu dowolnego późniejszego wykonania.
Przechowanie referencji do argumentu tymczasowego w strukturze,
która przeżyje wywołanie, może prowadzić do wskaźników lub
referencji wiszących. Ten sam problem dotyczy między innymi
\texttt{std::forward\_as\_tuple}.

\section{Tworzenie obiektów z dowolnych argumentów}
\label{sec:variadic-emplace}

\subsection{Własna fabryka obiektów}

\begin{lstlisting}
#include <memory>
#include <utility>

template<class T, class... Args>
std::unique_ptr<T> create(Args&&... args) {
    return std::make_unique<T>(
        std::forward<Args>(args)...
    );
}

struct Point {
    int x, y;
    Point(int a, int b) : x(a), y(b) {}
};

auto point = create<Point>(4, 7);
\end{lstlisting}

\texttt{std::make\_unique} już dostarcza taką funkcjonalność,
więc \texttt{create} ma tu charakter dydaktyczny. Pokazuje
jednak wzorzec używany w fabrykach oraz funkcjach
\texttt{emplace} i \texttt{make}.

\subsection{Tworzenie kontenera bez wymogu kopiowalności}

\begin{lstlisting}
#include <concepts>
#include <string>
#include <utility>
#include <vector>

template<class T, class... Args>
    requires (std::constructible_from<T, Args&&> && ...)
std::vector<T> make_vector(Args&&... args) {
    std::vector<T> result;
    result.reserve(sizeof...(Args));
    (result.emplace_back(std::forward<Args>(args)), ...);
    return result;
}

auto words = make_vector<std::string>("A", "BB", "CCC");
\end{lstlisting}

Operator przecinka daje tu jednoznaczną kolejność wywołań
\texttt{emplace\_back}. Funkcja przyjmuje także pusty pakiet
argumentów. W przeciwieństwie do budowania wektora z listy
inicjalizacyjnej pozwala łatwo przekazywać do kolejnych
\texttt{emplace\_back} argumenty przenoszalne, ale niekopiowalne.
Dla elementów konstruowanych z więcej niż jednego parametru
wymagałaby innego interfejsu grupowania argumentów.

\section{\texttt{std::tuple}, \texttt{std::apply} i \texttt{std::index\_sequence}}
\label{sec:variadic-tuple}

\subsection{Heterogeniczna krotka \texttt{std::tuple}}

\texttt{std::tuple<Ts...>} to kontener o strukturze opisanej
listą typów znaną w czasie kompilacji. Każda pozycja może
mieć inny typ:

\begin{lstlisting}
#include <tuple>
#include <string>

std::tuple<int, std::string, double> t{7, "Ada", 3.5};

int n = std::get<0>(t);
const std::string& name = std::get<1>(t);
\end{lstlisting}

W przeciwieństwie do \texttt{std::vector<T>}, krotka nie
przechowuje dowolnej liczby elementów jednego typu, lecz
\emph{ustaloną na etapie kompilacji} listę typów.

\subsection{Rozwinięcie krotki do argumentów funkcji}

Funkcja \texttt{std::apply} (C++17) wywołuje funkcję,
przekazując elementy obiektu \emph{tuple-like} jako osobne
argumenty:

\begin{lstlisting}
#include <iostream>
#include <tuple>

void apply_example() {
    std::apply([](const auto&... xs) {
        ((std::cout << xs << ' '), ...);
        std::cout << '\n';
    }, t);
}
\end{lstlisting}

\texttt{std::apply} oszczędza ręcznego konstruowania
sekwencji indeksów, ale warto wiedzieć, jak taki mechanizm
można zaimplementować.

\subsection{Pakiety indeksów: \texttt{std::index\_sequence}}

\begin{lstlisting}
#include <cstddef>
#include <iostream>
#include <tuple>
#include <utility>

template<class Tuple, std::size_t... I>
void print_tuple_impl(const Tuple& t,
                      std::index_sequence<I...>) {
    ((std::cout << std::get<I>(t) << ' '), ...);
    std::cout << '\n';
}

template<class... Ts>
void print_tuple(const std::tuple<Ts...>& t) {
    print_tuple_impl(t, std::index_sequence_for<Ts...>{});
}

void tuple_example() {
    print_tuple(std::make_tuple(1, "abc", 3.0));
}
\end{lstlisting}

\texttt{std::index\_sequence<0,1,2>} jest specjalizacją
szablonu przyjmującego pakiet indeksów całkowitych.
\texttt{std::index\_sequence\_for<Ts...>} generuje indeksy
$0,1,\ldots,n-1$ dla pakietu długości $n$.
Wzorzec \texttt{std::get<I>(t)...} rozwija wyrażenie
osobno dla każdego indeksu.

\subsection{Równoległe rozwijanie pakietów}

W jednym wzorcu można użyć kilku pakietów. Muszą mieć
\emph{tę samą długość}, jeśli są rozwijane równolegle:

\begin{lstlisting}
#include <tuple>
#include <type_traits>
#include <utility>

template<class A, class B>
struct zipped_types;

template<class... Ts, class... Us>
struct zipped_types<std::tuple<Ts...>, std::tuple<Us...>> {
    static_assert(sizeof...(Ts) == sizeof...(Us));
    using type = std::tuple<std::pair<Ts, Us>...>;
};

using Z = zipped_types<
    std::tuple<int, char>,
    std::tuple<double, bool>
>::type;

static_assert(std::is_same_v<
    Z, std::tuple<std::pair<int, double>,
                  std::pair<char, bool>>>);
\end{lstlisting}

Elementy \texttt{Ts} i \texttt{Us} są łączone pozycjami:
$\bigl(T_1,U_1\bigr),\ldots,\bigl(T_n,U_n\bigr)$.
Nie jest to iloczyn kartezjański pakietów.

\section{Dziedziczenie wariadyczne i \texttt{std::visit}}
\label{sec:variadic-overloaded}

Pakiet \texttt{Bases...} można rozwinąć na liście baz.
Z kolei zapis \texttt{using Bases::operator()...}
wprowadza ich operatory wywołania do jednej klasy.
W ten sposób otrzymujemy zestaw przeciążeń.

\begin{lstlisting}
#include <iostream>
#include <string>
#include <variant>

template<class... Fs>
struct overloaded : Fs... {
    using Fs::operator()...;
};

template<class... Fs>
overloaded(Fs...) -> overloaded<Fs...>; // C++17 deduction guide

int main() {
    std::variant<int, std::string> value = "hello";

    std::visit(overloaded{
        [](int n) { std::cout << "int: " << n << '\n'; },
        [](const std::string& s) {
            std::cout << "string: " << s << '\n';
        }
    }, value);
}
\end{lstlisting}

Jest to standardowy idiom C++17: każda lambda definiuje
osobny \texttt{operator()}, a dziedziczenie wariadyczne
łączy je w jeden obiekt wywoływalny. Gdy stosujemy ten wzorzec
w wielu miejscach, warto umieścić \texttt{overloaded} w
niewielkim nagłówku narzędziowym.

\section{Ograniczanie pakietów za pomocą \texttt{concepts} (C++20)}
\label{sec:variadic-concepts}

Nie każda operacja ma sens dla każdej możliwej kombinacji
typów. W C++20 można czytelnie określać wymagania dla
\emph{wszystkich} elementów pakietu.

\begin{lstlisting}
#include <concepts>
#include <type_traits>

template<class... Ts>
    requires (std::integral<Ts> && ...)
constexpr bool all_nonnegative(Ts... xs) {
    return ((xs >= 0) && ...);
}

static_assert(all_nonnegative(0, 2, 7));
static_assert(!all_nonnegative(0, -1, 7));

// all_nonnegative(3.14); // rejected by constraints
\end{lstlisting}

Możemy również ograniczać pakiet w odniesieniu do
konkretnego typu docelowego:

\begin{lstlisting}
#include <concepts>

template<class T, class... Args>
concept ConstructibleWith =
    std::constructible_from<T, Args...>;

static_assert(ConstructibleWith<std::string, const char*>);
static_assert(!ConstructibleWith<int, std::string>);
\end{lstlisting}

Pierwszy przykład stosuje fold nad \emph{warunkami dla pojedynczych
typów}; drugi sprawdza \emph{jedną własność całej sekwencji
typów}, tj. czy odpowiedni konstruktor przyjmuje argumenty
\texttt{Args...} w tej kolejności. Tych dwóch mechanizmów
nie należy utożsamiać.

\subsection{\texttt{if constexpr} i struktura pakietu}

\begin{lstlisting}
#include <iostream>

template<class... Ts>
void report(const Ts&... xs) {
    if constexpr (sizeof...(Ts) == 0) {
        std::cout << "empty\n";
    } else {
        ((std::cout << xs << ' '), ...);
        std::cout << '\n';
    }
}
\end{lstlisting}

\texttt{if constexpr} pozwala w zależności od długości pakietu
wybierać gałąź w czasie kompilacji. W dobrze sformułowanym
szablonie odrzucona gałąź nie jest instancjonowana dla danej
specjalizacji. To użyteczne przy przypadkach brzegowych,
które nie mają sensu dla pustego pakietu.

\section{Ograniczenia, pułapki i diagnostyka}
\label{sec:variadic-pulapki}

\subsection{Pakiet typów nie jest pakietem argumentów funkcji}

\begin{lstlisting}
template<class... Ts>
void f(Ts... args) {
    // sizeof...(Ts): number of types
    // sizeof...(args): number of function parameters
    // args[0];  // ERROR: args is not an array
}
\end{lstlisting}

Do wybierania argumentu na ustalonej pozycji można użyć
\texttt{std::tuple} i \texttt{std::get} albo wywoływać funkcję
przez \texttt{std::apply}. Sam pakiet \texttt{args...}
nie jest indeksowanym kontenerem.

\subsection{Brak zgodności typów po rozwinięciu}

\begin{lstlisting}
template<class... Ts>
auto add(Ts... xs) {
    return (0 + ... + xs);
}

// add(1, std::string{"x"}); // ERROR: unsuitable operator+
\end{lstlisting}

Pakiet nie omija systemu typów. Każde wyrażenie powstałe po
jego rozwinięciu musi być poprawne, podobnie jak zwykły kod
C++. Jeżeli szablon ma określone wymagania, najlepiej
zapisać je za pomocą \texttt{requires}.

\subsection{Niewłaściwe przekazywanie kategorii wartości}

\begin{lstlisting}
template<class F, class... Args>
void bad_call(F&& f, Args&&... args) {
    std::forward<F>(f)(args...); // loses rvalue categories
}

// Correct:
// std::forward<F>(f)(std::forward<Args>(args)...);
\end{lstlisting}

Dla typów niekopiowalnych taki błąd często ujawnia się
natychmiast podczas kompilacji. Dla typów kopiowalnych
może spowodować jedynie dodatkowe kopie albo zmianę
wybranego przeciążenia.

\subsection{Wnioskowanie typów i \texttt{\{...\}}}

\begin{lstlisting}
#include <initializer_list>

template<class... Args>
void foo(Args&&...) {}

void braced_example() {
    // foo({1, 2, 3}); // ERROR: cannot deduce Args
    foo(std::initializer_list<int>{1, 2, 3}); // OK
}
\end{lstlisting}

Sama lista ujęta w nawiasy klamrowe \texttt{\{...\}} nie ma
typu, który dowolny pakiet \texttt{Args\&\&...} potrafiłby
wydedukować. Potrzebny jest określony typ parametru,
np. \texttt{std::initializer\_list}\allowbreak\texttt{<int>}, lub jawnie
utworzony obiekt.

\subsection{Ograniczenia położenia pakietu}

W \emph{głównej definicji szablonu klasy} pakiet parametrów
szablonu musi być ostatni. Dla szablonów funkcji istnieją
bardziej elastyczne reguły, ale kolejne parametry muszą
być dedukowalne lub mieć argumenty domyślne, stosownie do
kontekstu:

\begin{lstlisting}
// ERROR for a primary class template:
// template<class... Ts, class U> struct invalid;

template<class U, class... Ts>
struct valid {};
\end{lstlisting}

\subsection{Skomplikowany kod zamiast prostego rozwiązania}

Szablon wariadyczny to narzędzie, a nie obowiązkowy styl.
Jeżeli argumenty mają \emph{jednolity typ i ich liczba
zmienia się w czasie wykonania}, lepszy bywa
\texttt{std::span<T>}, \texttt{std::vector<T>} lub para
iteratorów. Pakiet ma liczbę elementów i strukturę typów
ustalaną na etapie kompilacji; nie zastępuje kontenera.

\section{Variadic templates a funkcje wariadyczne w stylu C}

Historyczne funkcje z deklaracją \texttt{void f(int, ...)}
korzystają z mechanizmu argumentów wariadycznych języka C,
a nie z pakietu parametrów szablonu C++:

\begin{lstlisting}
#include <cstdio>

int main() {
    std::printf("%d %f\n", 7, 2.5); // C-style variadic call
}
\end{lstlisting}

W formacie C zgodność argumentów ze specyfikatorami
formatowania jest obowiązkiem programisty (niektóre
kompilatory potrafią wykrywać część pomyłek). W szablonach
wariadycznych typy argumentów pozostają dostępne
kompilatorowi, dzięki czemu można wymuszać ograniczenia
i stosować przeciążanie oraz \texttt{std::forward}.
Nie należy jednak wyciągać z tego wniosku, że dowolny
szablon automatycznie bezbłędnie formatuje argumenty.

\section{Pełny przykład: prosty rejestrator zdarzeń}
\label{sec:variadic-projekt}

Zastosujmy omówione mechanizmy do funkcji tworzącej tekstowy
komunikat z dowolnej liczby argumentów możliwych do
wypisania do strumienia. Zamiast ręcznie doklejać spacje
w wielu przeciążeniach, użyjemy fold nad operatorem przecinka.

\begin{lstlisting}
#include <iostream>
#include <sstream>
#include <string>

// Each argument must support insertion into std::ostream.
template<class... Args>
std::string make_message(const Args&... args) {
    std::ostringstream out;
    ((out << args), ...);
    return out.str();
}

template<class... Args>
void log_line(const Args&... args) {
    std::cout << "[LOG] " << make_message(args...) << '\n';
}

int main() {
    log_line("count=", 3, ", ratio=", 0.75);
    log_line("finished");
    log_line();
}
\end{lstlisting}

\paragraph{Analiza.}
Dla argumentów \texttt{("count=",3,", ratio=",0.75)}
pakiet zawiera cztery elementy. Wyrażenie
\texttt{((out << args), ...)} wstawia je kolejno do strumienia,
a \texttt{make\_message} zwraca powstały napis. Dla
pustego pakietu fold z operatorem przecinka jest poprawny,
a wynikowy napis jest pusty. W wersji produkcyjnej należałoby
rozważyć wymagania dotyczące formatowania, zachowanie w
obliczu wyjątków, współbieżność oraz ewentualne skorzystanie
z narzędzia formatującego wspieranego przez używany kompilator.

\section{Zestawienie: kiedy używać poszczególnych mechanizmów?}

\begin{center}
\small
\begin{tabular}{|p{0.40\textwidth}|p{0.46\textwidth}|}
\hline
\textbf{Cel} & \textbf{Typowe rozwiązanie} \\
\hline
Dowolna liczba argumentów o różnych typach & \texttt{template<class... Ts>} i \texttt{Ts...} \\
\hline
Wykonanie tej samej operacji dla każdego argumentu & Fold z operatorem \texttt{,} (C++17) \\
\hline
Suma, iloczyn lub wspólny warunek & Fold z \texttt{+}, \texttt{*}, \texttt{\&\&} lub \texttt{||} \\
\hline
Przekazanie argumentów dalej z zachowaniem kategorii & \texttt{Args\&\&...}, \texttt{std::forward<Args>(args)...} \\
\hline
Zapis heterogenicznych wartości & \texttt{std::tuple<Ts...>} \\
\hline
Wywołanie funkcji elementami krotki & \texttt{std::apply} (C++17) \\
\hline
Rozwinięcie krotki przy pomocy indeksów & \texttt{std::index\_sequence} \\
\hline
Lista typów do obliczeń podczas kompilacji & Wariadyczny szablon klasy i specjalizacje \\
\hline
Ograniczenie do odpowiednich typów & Fold z koncepcjami i \texttt{requires} (C++20) \\
\hline
Zmienna długość listy jednolitych danych w czasie wykonania & \texttt{std::vector} lub \texttt{std::span}, nie pakiet \\
\hline
\end{tabular}
\end{center}

\section{Podsumowanie i zadania do samodzielnego sprawdzenia}

\begin{enumerate}
    \item \texttt{class... Ts} deklaruje \emph{pakiet typów},
    a \texttt{Ts... xs} wprowadza \emph{pakiet parametrów funkcji}.
    \item \texttt{sizeof...(Ts)} oblicza liczbę elementów pakietu
    podczas kompilacji.
    \item \texttt{pattern...} rozwija cały \emph{wzorzec},
    nie tylko pojedynczy identyfikator.
    \item Fold expressions (od C++17) zastępują wiele prostych
    rekurencyjnych implementacji z C++11.
    \item Dla pustego pakietu istotne są element neutralny i
    reguły operatora; nie każda forma fold jest poprawna.
    \item W wrapperach zachowuj kategorię wartości argumentów,
    używając referencji przekazujących i \texttt{std::forward}.
    \item \texttt{std::tuple}, \texttt{std::apply} i
    \texttt{std::index\_sequence} współpracują z pakietami,
    umożliwiając operacje na heterogenicznych danych.
    \item W C++20 używaj \texttt{concepts}/\texttt{requires},
    aby błędne wywołania odrzucać na granicy interfejsu.
\end{enumerate}

\paragraph{Pytania kontrolne.}
Spróbuj przewidzieć wyniki poniższych wyrażeń bez uruchamiania
programu. Następnie zweryfikuj je za pomocą \texttt{static\_assert}:

\begin{lstlisting}
// static_assert(sizeof...(int) == 1); // ERROR: int is not a pack

// For the functions defined above:
static_assert(product() == 1);
static_assert(left_sub(10, 3) == 87);
static_assert(right_sub(10, 3) == 107);
static_assert(count_integral_v<int, double, char> == 2);
\end{lstlisting}

Pierwsza linia jest \emph{celowo niepoprawna}: operator
\texttt{sizeof...} wymaga nazwy pakietu, a nie zwykłego typu.

% Wykorzystane punkty odniesienia (bez kopiowania tresci):
% https://en.cppreference.com/w/cpp/language/parameter_pack
% https://en.cppreference.com/w/cpp/language/fold
% https://en.cppreference.com/w/cpp/utility/forward
% https://en.cppreference.com/w/cpp/utility/apply
% https://en.cppreference.com/w/cpp/utility/integer_sequence
